\documentclass[12pt,a4paper,oneside]{article}
  %\documentclass{amsart}
  \usepackage{authblk} %for writing authors
  \usepackage{url}
 % \usepackage{pgfplots}
  \usepackage{amsthm}
  \usepackage{amsmath}
  \usepackage{amsfonts}
  \usepackage{amssymb}
  \usepackage{amscd}
  \usepackage{authblk}
  \usepackage{bm}
  \usepackage[mathscr]{eucal}%for \mathscr font
  \usepackage{mathtools} % for 'bsmallmatrix' environment
  \usepackage{enumitem}
  \usepackage{physics}%for using Dirac's bra-ket notation, code: $\bra{\Psi}\ket{\Psi}$, $\expval{A}{\Psi}$, $\braket{x}{y}$
  \usepackage{color}
 % \usepackage[affil-it]{authblk} % author affiliation etc.
  \usepackage{fancyhdr}
  \usepackage{hyperref}
  %\usepackage[colorlinks=true, linkcolor=blue]{hyperref}
  %\usepackage{showkeys}
  \usepackage{marginnote}%for comments in margin, use \marginpar{text}
  \usepackage[top=2.5cm, bottom=2.7cm,left=2.5cm, right=2.5cm, marginparwidth=1.8cm]{geometry}
  %\usepackage[top=2.5cm, bottom=2.7cm,left=3.5cm, right=2.5cm]{geometry}
  \usepackage{framed}% use \begin{frame} text \end{frame} to put it in a framed box
  \usepackage[noabbrev]{cleveref} %for numbering in align* use \cref{...} for refering
  \usepackage[T1]{fontenc} %for writing multiple authors with multiple affiliations
 \usepackage[utf8]{inputenc}%for writing multiple authors
 \usepackage{blkarray}%for block matrices with subspaces marked for example see below


\usepackage{relsize} % example: \mathlarger{\mathlarger{\Pi}}
                     % documentation is available  at https://ctan.math.washington.edu/tex-archive/macros/latex/contrib/relsize/relsize-doc.pdf 
  
  %\usepackage{showkeys}
  
 
 %Code for numbering equations sectionwise
  \numberwithin{equation}{section}
 
 \allowdisplaybreaks[1] 
  
\pagestyle{plain}


  \theoremstyle{definition}  %Title and number in bold, body in normal font.
   \newtheorem{defn}{Definition}[section]
   \newtheorem{eg}[defn]{Example}
   \newtheorem{egs}[defn]{Examples}
   \newtheorem*{prblm}{Problem}
   \newtheorem*{prblms}{Problems}
   \newtheorem*{ex}{Exercise}
   \newtheorem{exs}[defn]{Exercises}
   \newtheorem*{note}{Note}
   \newtheorem*{fact}{Fact}
   \newtheorem{rmk}[defn]{Remark}
   \newtheorem{rmks}[defn]{Remarks}

  \theoremstyle{plain}  %n Title and number in bold, body in italic (default).
   \newtheorem{thm}[defn]{Theorem}
   \newtheorem{lem}[defn]{Lemma}
   \newtheorem{prop}[defn]{Proposition}
   \newtheorem{cor}[defn]{Corollary}
   \newtheorem*{conj}{Conjecture}
   \newtheorem*{ackn}{Acknowledgment}
   \newtheorem*{qn}{Question}

  \theoremstyle{remark} %Title and number in italic, body in normal font.
   \newtheorem*{notn}{Notation}
   \newtheorem*{summ}{Summary}
   \newtheorem*{concl}{Conclusion}
   \newtheorem*{claim}{Claim}
   \newtheorem*{pf}{Proof}
   

      \newcommand{\me}{\mathrm{e}}
   \newcommand{\B}[1]{\mathscr{B}({#1})}
   \newcommand{\Bo}[1]{\mathscr{B}_1({#1})}
   \newcommand{\Bt}[1]{\mathscr{B}_2({#1})}
   \newcommand{\Bp}[1]{\mathscr{B}_p({#1})}
   \newcommand{\BRH}{\mathscr{B}_{\mathbb{R}}(\mathcal{H})}
   \newcommand{\SH}[1]{\mathscr{S}({#1})}
    \newcommand{\cv}[1]{\mathbf{#1}}
   \newcommand{\ez}{e(\cv{z})}
   \newcommand{\eu}{e(\cv{u})}
   %\newcommand{\ev}{e(\cv{v})} 
   \newcommand{\ew}{e(\cv{w})}
   \newcommand{\wz}{W(\cv{z})}
   \newcommand{\wu}{W(\cv{u})}
   \newcommand{\wv}{W(\cv{v})}
   \newcommand{\bell}{\bm{\ell}}
   \newcommand{\bk}{\mathbf{k}}
   \newcommand{\br}{\mathbf{r}}
   \newcommand{\bs}{\mathbf{s}}
   \newcommand{\bt}{\mathbf{t}}
   \newcommand{\bu}{\mathbf{u}}
   \newcommand{\bv}{\mathbf{v}}
   \newcommand{\q}{\mathbf{q}}
   \newcommand{\m}{\mathbf{m}}
   \newcommand{\w}{\mathbf{w}}
   \newcommand{\x}{\mathbf{x}}
   \newcommand{\y}{\mathbf{y}}
   \newcommand{\z}{\mathbf{z}}
   \newcommand{\F}{\mathscr{F}}
   \newcommand{\Fd}{\mathscr{F}_{d}}
   \newcommand{\COH}{\mathscr{C}(\mathcal{H})}
   \newcommand{\CE}{\mathcal{E}}
   \newcommand{\CH}{\mathcal{H}}
   \newcommand{\CK}{\mathcal{K}}
   \newcommand{\GH}{\Gamma(\mathcal{H})}
   \newcommand{\GC}{\Gamma(\mathbb{C})}
   \newcommand{\GCk}{\Gamma(\mathbb{C}^{k})}
   \newcommand{\GCn}{\Gamma(\mathbb{C}^{n})}
   \newcommand{\GCd}{\Gamma(\mathbb{C}^{d})}
   \newcommand{\GOL}{\Gamma_{0}(L)}
   \newcommand{\BC}{\mathbb{C}}
   \newcommand{\BR}{\mathbb{R}}
   \newcommand{\BF}{\mathbb{F}}
   \newcommand{\BN}{\mathbb{N}}
   \newcommand{\BZ}{\mathbb{Z}}

   
\usepackage{centernot}   
\newcommand{\nll}{\centernot{\ll}}

 \newcommand{\numberthis}{\refstepcounter{equation}\tag{\theequation}}  
\makeatletter
\newcommand\footnoteref[1]{\protected@xdef\@thefnmark{\ref{#1}}\@footnotemark}
\makeatother 

\renewcommand\Authfont{\scshape\small}
\renewcommand\Affilfont{\itshape\small}
\setlength{\affilsep}{1em}

\newcommand{\keywords}[1]{{\footnotesize\hspace{0.68cm}{\textit{Keywords}: }#1\par
  \vskip.7\baselineskip}}
\newcommand{\emailaddress}[1]{\newline{\sf#1}}


\renewcommand{\colon}{\nobreak\mskip2mu\mathpunct{}\nonscript
  \mkern-\thinmuskip{:}\mskip6muplus1mu\relax}
\let\tr\relax 
\DeclareMathOperator{\tr}{Tr}
\DeclareMathOperator{\diag}{Diag}
\let\ker\relax 
\DeclareMathOperator{\ker}{Ker}
\DeclareMathOperator{\ran}{Ran}
\DeclareMathOperator{\re}{Re}
\DeclareMathOperator{\im}{Im}
\DeclareMathOperator{\supp}{Supp}
\DeclareMathOperator*{\slim}{s-lim}
\DeclareMathOperator{\spn}{span}
\DeclareMathOperator{\variance}{Var}
\DeclareMathOperator{\cov}{Cov}
\fontsize{12}{14}



\title{Petz-R\'enyi Relative Entropy of Quantum States  Reduces to R\'enyi Relative Entropy of Classical Probability Measures}
\author[1]{George Androulakis}
\author[2]{Tiju Cherian John}


\affil[1,2]{University of South Carolina, Columbia, South Carolina, USA \emailaddress{giorgis@math.sc.edu, tijucherian@fulbrightmail.org}}



\begin{document}
\maketitle
\tableofcontents
\newpage
 \begin{abstract}
Let $\mathcal{K}$ be a complex separable Hilbert space of finite or infinite dimensions. We prove that the Petz-R\'enyi relative entropy of any two quantum states $\rho$ and $\sigma$ on $\mathcal{K}$, denoted by $D_{\alpha}(\rho||\sigma)$ reduces to the R\'enyi relative entropy (divergence) of two classical probability measures $P$ and $Q$ obtained from $\rho$ and $\sigma$. This leads to a number of new results in the infinite dimensions,  and new proofs for some known results in the finite dimensional setting. Our method provides a general framework for proving a quantum counterpart of any result about the classical R\'enyi divergence involving only two probability distributions.  Furthermore, we construct an example to show that the information theoretic definition of the von Neumann relative entropy is different from Araki's definition of relative entropy. This disproves a recent attempt in the literature to prove  such a result. 
All the results proved here  are valid in both finite and infinite dimensions and hence can be applied to continuous variable systems as well.\\ 
   \textbf{Keywords:}  Relative divergence, Petz-R\'enyi relative entropy, R\'enyi divergence, von Neumann relative entropy, Araki relative entropy, quantum Pinsker's inequality, continuous variable quantum information\\
 \textbf{2020 Mathematics Subject classification:}
 Primary 81P17; Secondary 81P99.
 \end{abstract}
 
 
\section{Introduction}\label{sec:introduction}
Following the pioneering work of  Shannon on information theory, Kullback and Leibler defined  a divergence of two probability measures $P$ and $Q$ \cite{kullback-leibler-1951} \[D_{KL}(P||Q)= \sum_{x} P(x) \log \frac{P(x)}{Q(x)} \text{ with the convention } 0 \cdot \infty =0.\] This would later come to be known as the Kullback-Leibler divergence of probability measures. Later,  R\'enyi generalised the notion of Kullback-Leibler divergence \cite{renyi1961}. For $0< \alpha \not = 1$, R\'enyi  defined
\[
I_\alpha (P||Q)= \frac{1}{\alpha -1} \log \sum_x P(x)^\alpha Q(x)^{1-\alpha}.
\] 
Quantizing these developments Umegaki in \cite{umegaki-1962} extended these notions to the setting of von Neumann algebras. If
$\mathcal{M}$ be a von Neumann algebra with a faithful normal trace $\tau$, and $\rho, \sigma$ are elements of the noncommutative  $L^1(\mathcal{M},\tau)$ space, Umegaki defined 
\[
I(\rho||\sigma)= \left\{
\begin{array}{ll}
\tau (\rho \log \rho - \rho \log \sigma), &
\text{ if } \text{supp}\, \rho \subseteq \text{supp}\, \sigma;\\
\infty, & \text{ otherwise}.
\end{array}\right.
\] Araki defined the relative entropy between two  states of a von Neumann algebra using the theory of relative modular operators in
 \cite{araki1976relative} and \cite{araki1977relative} (See Appendix \ref{appendix:relative-modular} and Section \ref{sec:counter-eg}). Similar to the generalisation  R\'enyi obtained for the Kullback-Leibler divergence, Petz and Ohya  \cite{petz1986quasientropy, petz-ohya-1993} generalised the notion of Araki's relative entropy to what is known today as Petz-R\'enyi $\alpha$-relative entropy. The Petz-R\'enyi $\alpha$-relative entropy of two density operators (positive trace class operators with unit trace, also known as states) $\rho$ and $\sigma$ is defined as \[D_{\alpha} (\rho \lvert\lvert \sigma) = \frac{1}{\alpha-1}\log \tr \rho^{\alpha/2}\sigma^{1-\alpha}\rho^{\alpha/2}, \alpha \in (0,1)\cup (1,\infty). \] Datta in \cite{Datta2009} studied an entropic quantity called $D_{\operatorname{min}}(\rho||\sigma)$ in the finite dimensional setting, and verified that   $D_{\operatorname{min}}(\rho||\sigma) = \lim_{\alpha\rightarrow0}D_{\alpha}(\rho||\sigma)$  \cite[equation (10)]{Datta2009}. Recently, Seshadreeshan, Lami and Wilde in \cite{Ses-Lam-Wil-2018} and Parthasarathy in \cite{Par2021a} studied the Petz-R\'enyi relative entropy in the infinite dimesnional setting, for the particular case of gaussian states. A related entropic quantity called sandwiched R\'enyi relative entropy gained prominence in the recent years due to several nice properties satisfied by them,  \cite{muller-tomamichel-etal-2013,frank-leib-2013, wilde-winter-yang-2014, Datta_Leditzky_2014, Par2021b}.  It should be noted that all of these entropic quantities are important in quantum information theory in general and quantum hypothesis testing in particular. One may refer to \cite{petz-hiai-1991} and the other articles cited above to know more about different uses of relative entropy.
 
 Our aim in this article is to prove that the Petz-R\'enyi relative entropy enjoys several important properties  even in the infinite dimensional setting which is not well understood in the literature.
Some important results obtained in this article are: \begin{enumerate}
    \item \textit{A new formula:} For $\alpha\in (0,1)\cup (1,\infty)$, a formula to compute the Petz-R\'enyi $\alpha$-relative entropy between states $\rho$ and $\sigma$, in terms of certain  classical probability measures on the positive integer lattice of the plane $\BR^2$ (Theorems \ref{thm:the-formula-general}, and \ref{thm:quantum-reduces-to-classical}).
    
    \item \textit{Positivity and Monotonicity:} $D_{\alpha}(\rho||\sigma)\geq 0$ and  is  nondecreasing in $\alpha$  (Theorems \ref{thm:3-veh-quantum} and \ref{thm:8-veh-quantum}).
    
    \item \textit{Limiting Cases:} Formulae for limits of $D_{\alpha}(\rho||\sigma)$ as $\alpha\rightarrow 0, 1,$ and $\infty$. The von Neumann relative entropy $D(\rho||\sigma)$ is recovered as  \[\lim_{\alpha \uparrow 1} D_{\alpha}(\rho||\sigma) = D(\rho||\sigma)\](Theorems \ref{thm:4-veh-quantum}, \ref{thm:6-veh-quantum} anf \ref{thm:von-Neumann-entropy}).
    
   \item Necessary and sufficient conditions for $\lim_{\alpha\downarrow0}D_{\alpha}(\rho||\sigma) =0$ (Theorem \ref{thm:23-veh-quantum}).
   
    \item Necessary and sufficient conditions for $D_{\alpha}(\rho||\sigma) =\infty$ for some $\alpha\in [0,1)$ (Theorem \ref{thm:24-veh-quantum}).
  
    \item \textit{A counter example:} Information theoretic definition of the von Neumann relative entropy is different from Araki's definition of relative entropy. This disproves a recent attempt to prove  such a result in   \cite[Theorem 20]{luczak2019relative} (Example \ref{eg:counter-example}, Theorem \ref{thm:counter-example}).
    
    \item \textit{Continuity:}  $D_\alpha (\rho||\sigma)$  is continuous in $\alpha$ on \[\mathcal{A} = \{\alpha\in [0,\infty]\,:\, 0\leq \alpha\leq 1\textnormal{ or } D_{\alpha}(\rho||\sigma)<\infty\}\]
    (Theorem \ref{thm:7-veh-quantum}).
    
    \item \textit{Quantum Pinsker Type inequality:} For $\alpha\in (0,1]$ \begin{align*}
  \frac{\alpha}{2}\norm{\rho-\sigma}_2^4\leq   D_{\alpha}(\rho||\sigma),
 \end{align*}
 where $\norm{\cdot}_2$ denotes the Hilbert-Schmidt norm (Theorem \ref{thm:quantum-pinsker}).
    
    \item \textit{Skew symmetry:} for $0<\alpha<1$, \[D_\alpha(\rho||\sigma) = \frac{\alpha}{1-\alpha}D_{1-\alpha}(\sigma||\rho)\]
    (Proposition \ref{prop:8-veh-quantum}).
    
    \item \textit{Concavity:} The function $(1-\alpha)D_{\alpha}(\rho||\sigma)$ is concave in $\alpha$ (Proposition \ref{cor:2-veh-quantum}).
    \end{enumerate}
 
 \section{Preliminaries}\label{sec:preliminaries}


\subsection{A Criteria for Finite Trace}Let us recall the definition of the trace of a positive operator, as this definition have an important role in our analysis. We refer to Section 11 in Chapter 2 of \cite{vN2018} for this definition of the trace. Let $\CK$ denote an arbitrary but fixed complex separable Hilbert space for this section. In what follows, a selfadjoint operator $T$ (not necessarily bounded) is said to be positive  if \begin{equation}
    \mel{x}{T}{x}\geq 0,\quad \forall x\in D(T),
\end{equation}  (See Definition 9 of  Chapter 2 (page 62) of \cite{vN2018}). A positive operator defined in this way has spectrum on the positive half of the real line. 
\begin{defn}\label{defn:trace} 
 Let $T$ be a positive operator defined on a dense subspace $D(T)\subseteq \CK$. Let $\{u_j\} \subseteq D(T)$ be any orthonormal basis (onb) of $\CK$ (such an onb always exists by applying Gram-Schmidt process to a countable dense subset of $D(T)$). Define the trace of $T$, denoted by $\tr T$, by \begin{align}
     \tr T = \sum\limits_j \mel{u_j}{T}{u_j} \in [0, \infty].
 \end{align} 
 Trace of $T$ thus defined is independent of the choice of the onb $\{u_j\}$ \cite{vN2018}. A positive operator T is said to be a trace class operator if $\tr T <\infty$.
\end{defn}

The following lemma shows that a positive trace class operator (as in the previous definition) extends as a bounded operator on $\CK$.
\begin{lem}\label{sec:gaussian-states-lem}
 Let $D(T)\subseteq \CK$ be a dense subspace 
 and let $T:D(T)\rightarrow \CK$ be a positive operator. Assume that there exists an orthonormal basis $\{u_k\}\subseteq D(T)$ such that \begin{align*}
     \sum\limits_k \mel{u_k}{T}{u_k} = M<\infty,
 \end{align*} i.e., $T$ is a trace class operator in the sense of Definition \ref{defn:trace},
 then $T$ extends to a bounded  operator on $\CK$.
\end{lem}
\begin{proof}
It is enough to prove that $T$ is bounded. Let $h\in D(T)$ be a finite linear combination of $u_k$'s, i.e., $h = \sum_k \braket{u_k}{h}\ket{u_k}$, where $\braket{u_k}{h} = 0$ for all but finitely many $k$'s. Then \begin{align*}
    Th = \sum_k \braket{u_k}{h}T\ket{u_k}.
\end{align*}
Let $\xi^T$ denote the spectral measure associated with $T$, $\sqrt{T}$ exists as a positive operator and then we have \begin{align*}
    D(\sqrt{T}) &= \{u\in \CK: \int_{\BR}\abs{x}\mel{u}{\xi^T(\dd x)}{u}<\infty\}\\&
    \supseteq \{u\in \CK: \int_{\BR}\abs{x}^2\mel{u}{\xi^T(\dd x)}{u}<\infty\} = D(T).
\end{align*}
 Therefore, by applying Cauchy-Schwartz inequality twice we have
\begin{align*}
  0\leq  \mel{h}{T}{h} &= \braket{\sum \limits_j \braket{u_j}{h}\ket{u_j}}{\sum \limits_k \braket{u_k}{h}T\ket{u_k}}\\ 
    & = \sum\limits_{j,k}\overline{\braket{u_j}{h}}\braket{u_k}{h}\mel{u_j}{T}{u_k}\\
      & \leq \sum\limits_{j,k}\abs{\overline{\braket{u_j}{h}}\braket{u_k}{h}\mel{u_j}{T}{u_k}}\\
    &\leq  \sum\limits_{j,k}\abs{\overline{\braket{u_j}{h}}\braket{u_k}{h}\norm{\sqrt{T}u_j}\norm{\sqrt{T}u_k}}\\
    & = \left(\sum\limits_{k}\abs{\braket{u_k}{h}}\norm{\sqrt{T}u_k}\right)^2\\
    & \leq \left(\sum\limits_{k}\abs{\braket{u_k}{h}}^2\right)\left(\sum\limits_{k}\norm{\sqrt{T}u_k}^2\right)\\
    &= M \norm{h}^2, \quad \forall h\in D(T).
\end{align*}
Hence $T$ defined on $D(T)$ extends to a bounded operator on $\CK$.
\end{proof}
In order to study the Petz-R\'enyi $\alpha$-relative entropy (see Definition \ref{defn:alpha-relative-entropy}) of two states, $\rho$ and $\sigma$    we need to understand whether an operator of the form $\rho^{\alpha/2}\sigma^{1-\alpha}\rho^{\alpha/2}$ is trace class or not. Theorem \ref{sec:gaussian-states-prop} in this section provides an equivalent criteria to deal with this situation. We do not expect Theorem \ref{sec:gaussian-states-prop} to be new but we could not find a suitable reference to cite the same. First we recall the definition of the support of a positive operator.
\begin{defn}\label{defn:support}
 Define the support of a positive operator $\tau$, denoted by $\supp\tau$ as \[\supp\tau = (\ker\tau)^{\perp} = \overline{\ran}\tau.\]
\end{defn}
\begin{defn}\label{sec:states-observables-1}\label{defn:moments}Let $\CK$ be a complex separable Hilbert space and $\CK_1\subseteq \CK$ be a closed subspace. Let
  $X$ be a (possibly unbounded) selfadjoint operator defined on $D(X)\subseteq\CK_1\subseteq \CK$  with spectral measure $\xi^{X}$  and let $\tau$ be a positive compact operator on   $\CK$ with $\supp \tau \subseteq \CK_1$. Define a (possibly infinite) positive measure $\mu^{\tau,X}$ on $(\mathbb{R}, \mathcal{B}_{\mathbb{R}})$, where $\mathcal{B}_{\mathbb{R}}$ is the Borel $\sigma$-algebra of $\mathbb{R}$, by  \[\mu^{\tau,X}(E):= \tr \xi^{X}(E)\tau,\quad \forall E \in \mathcal{B}_{\mathbb{R}}.\] Then $\mu^{\tau,X}$ is called the \emph{distribution of  $X$ with respect to $\tau$}.
\end{defn}
\begin{rmk} \label{rmk:distribution}
If $\tau$ is a positive compact operator on $\CK$  with $\supp \tau \subseteq \CK_1$, then by the spectral theorem, 
\begin{equation}
\label{eq:15}
\tau = \sum_{i}p_{i}\ketbra{u_i}{u_{i}},
\end{equation}
 where $p_{i}\geq 0$ and  $\{u_i\}$ is an  orthonormal basis in $\CK$, $u_i\in \CK_1$ when $p_i\neq 0$. In this case, \begin{equation}
\label{eq:9}
\mu^{\tau,X}(E) = \sum\limits_{\{i\,:\, p_i\neq 0\}}^{}p_i \left\langle u_i|\xi^X(E)|u_i \right\rangle,
 \end{equation} 
 for any Borel set $E \subseteq \mathbb{R}$.  Thus $\mu^{\tau,X}$ is supported inside the spectrum of $X$. Furthermore, $\mu^{\tau,X}$ is a $\sigma$-finite measure.
\end{rmk}
\begin{thm}\label{sec:gaussian-states-prop}
   Let $\tau$ be a
  positive compact operator on $\CK$ with spectral decomposition 
\begin{align*}
\tau = \sum\limits_{i\in\mathcal{I}}^{}p_i\ketbra{u_i},
\end{align*} where $p_i\geq 0$ and $\{u_i\}_{i\in\mathcal{I}}$ an orthonormal basis of
$\CK$. Let $X$ be a positive operator defined on a dense subspace $D(X)$ of a closed subspace $\mathcal{K}_1$ of $\mathcal{K}$. 
Assume that  \begin{align}\label{item:8}
    u_i\in D(X), \text{ for all } i\in \mathcal{I} \text{ such that } p_i\neq 0.
\end{align}Then the operator $\tau^{1/2}X\tau^{1/2}$ initially defined on the dense subspace $\spn\{u_i\}_{i\in \mathcal{I}}$ by \begin{align*}
   \tau^{1/2}X\tau^{1/2}u_i= \begin{cases}
     \sqrt{p_i} \tau^{1/2}Xu_i,  &  \text{ when } p_i\neq 0\\
     0, &\text{ when } p_i = 0
      \end{cases}
\end{align*} extends as a  positive selfadjoint operator on a dense domain in  $\CK$. Furthermore,
\begin{align}
    \label{item:7} \tr \tau^{1/2}X\tau^{1/2} = \int_{0}^{\infty}x \mu^{\tau,X}(\dd x),
\end{align}
where $\mu^{\tau,X}$ is the distribution of $X$ with respect to $\tau$ (see Definition \ref{sec:states-observables-1}).
\end{thm}
\begin{proof}
Note that $\spn\{u_i\}_{i\in \mathcal{I}}$ is a dense subspace of $\CK$. Clearly, $\tau^{1/2}X\tau^{1/2}$ defined as in the statement of the Theorem is a symmetric operator. For any $x\in \spn\{u_i\}_{i\in \mathcal{I}}$, we have \[\mel{x}{\tau^{1/2}X\tau^{1/2}}{x} = \mel{\tau^{1/2}x}{X}{\tau^{1/2}x} \geq 0.\]
Now the fact that  $\tau^{1/2}X\tau^{1/2}$ defined on $\spn\{u_i\}_{i\in \mathcal{I}}$ extends as a positive selfadjoint operator is a consequence of the Friedrichs extension (\cite[Theorem 5.1.13]{pedersen2012analysis}, also \cite[Theorem X.23]{ReS75ii}). We fix the Friedrichs extension  of $\tau^{1/2}X\tau^{1/2}$ and use same notation for the extended operator.  

Now we prove (\ref{item:7}). Since $\tau^{1/2}X\tau^{1/2}$ is a positive operator, we have
\begin{align*} \tr \tau^{1/2}X\tau^{1/2}&=
\sum\limits_{i\in \mathcal{I}}^{}\mel{u_i}{\tau^{1/2}X\tau^{1/2}}{u_i}\\ &=
                                                        \sum\limits_{\{i|p_i\neq 0\}}^{}p_i\mel{u_i}{X}{u_i}\\
  & =\sum\limits_{\{i|p_i\neq 0\}}^{}p_i \int\limits_{0}^{\infty}x\mel{u_i}{\xi^X(\dd
    x)}{u_i}\numberthis \label{eq:20}\\
  & = \int\limits_{0}^{\infty}x \sum\limits_{\{i|p_i\neq 0\}}^{}p_i\mel{u_i}{\xi^X(\dd
    x)}{u_i}\numberthis \label{eq:21.1}\\
& = \int\limits_{0}^{\infty}x \sum\limits_{\{i|p_i\neq 0\}}^{}p_i\mel{u_i}{\xi^X(\dd
    x)}{u_i}\\    
  &=\int_{0}^{\infty}x \mu^{\tau,X}(\dd x),
\end{align*} where (\ref{eq:20}) follows from the spectral theorem, (\ref{eq:21.1}) follows because all the terms are positive and we used (\ref{eq:9}) in the last line above. 
\end{proof}
\begin{rmk}
Since
$T=\tau^{1/2}X\tau^{1/2}$ is a positive operator, Lemma
\ref{sec:gaussian-states-lem} shows that $\tau^{1/2}X\tau^{1/2}$ is a
trace class operator if and only if $\int_{0}^{\infty}x \mu^{\tau,X}(\dd x)<\infty$.
\end{rmk}
 \section{Petz-R\'enyi Relative Entropy of States}
 \begin{defn}\label{defn:alpha-relative-entropy}
Let $\rho$ and $\sigma$ be any two states on a Hilbert space $\CK$. For $\alpha \in (0,1)\cup (1,\infty)$, the Petz-R\'enyi $\alpha$-relative entropy of $\rho$ given $\sigma$ is 
\begin{equation}\label{eq:1.1}
  D_{\alpha} (\rho \lvert\lvert \sigma) = \frac{1}{\alpha-1}\log \tr \rho^{\alpha/2}\sigma^{1-\alpha}\rho^{\alpha/2}, 
\end{equation}
whenever $ \rho^{\alpha/2}\sigma^{(1-\alpha)}\rho^{\alpha/2} $ is a positive selfadjoint operator defined on a dense subspace operator on $\CK$. Otherwise, $D_{\alpha} (\rho \lvert\lvert \sigma) =\infty$.
\end{defn}
\begin{rmks}\begin{enumerate}
    \item 
When $\alpha>1$, we note that $\sigma^{(1-\alpha)}$ is defined as the  pseudoinverse of $\sigma$ raised to the power $(\alpha-1)$.  If the spectral decomposition of $\sigma$ is $\sum_{j}\lambda_j\ketbra{f_j}$ with $\lambda_j>0$, then \begin{align}\label{eq:pseudo-inverse}
    \sigma^{(1-\alpha)} = \sum_{j}\lambda_j^{(1-\alpha)}\ketbra{f_j}.
\end{align}
It may be noted from the spectral theorem (Theorem 12.4 in \cite{Par12}) that the pseudo inverse as defined above is a selfadjoint operator (not necessarily bounded) because its spectral measure is supported on the  real line. Furthermore, by (\ref{eq:pseudo-inverse}), we also have $\sigma^{(1-\alpha)}$ is a positive operator.

\item Let $\norm{A}_p$ denote the Schatten $p$-norm of an operator $A$.  When $0<\alpha<1$, we know that $\rho^{\alpha}$ and $\sigma^{(1-\alpha)}$ are bounded operators. In this case we have,
\[ \tr \rho^{\alpha/2}\sigma^{1-\alpha}\rho^{\alpha/2} = \tr \rho^{\alpha}\sigma^{1-\alpha}\\
\leq \norm{\rho^{\alpha}}_{\frac{1}{\alpha}}\norm{\sigma^{1-\alpha}}_{\frac{1}{1-\alpha}}=1.\]

Hence $0\leq D_{\alpha}(\rho||\sigma)$ for $0<\alpha<1$.
\end{enumerate}
\end{rmks}
\subsection{Main Results}
\begin{thm}\label{thm:the-formula-general}
Let $\CK$ be a complex Hilbert space with  $\dim \CK = \abs{\mathcal{I}}$, where $\mathcal{I}\subseteq \BN$ may be a finite or infinite set. Let $\rho$ and $\sigma$ be states on $\CK$  with spectral decomposition \begin{align}\label{eq:spectral-rho-sigma}\begin{split}
    \rho &= \sum_{i\in \mathcal{I}}r_i \ketbra{u_i}, \quad r_i\geq 0,\quad \sum_i r_i = 1,\quad \{u_i\}_{i} \text{ is an orthonormal basis};\\
    \sigma &= \sum_{j\in \mathcal{I}}s_j \ketbra{v_j}, \quad s_j\geq0, \quad \sum_j s_j = 1, \quad \{v_j\}_j\text{ is an orthonormal basis}.
    \end{split}
\end{align}
 Let $\alpha\in (0,1)\cup(1,\infty)$. Then the Petz-R\'enyi $\alpha$-relative entropy of $\rho$ and $\sigma$ is finite if and only if the following two conditions are satisfied: \begin{enumerate}
    \item\label{item:1-general} $u_i\in D(\sigma^{(1-\alpha)})$, $\forall i $ such that $r_i\neq 0$, which is same as \begin{equation}\label{eq:condition-1}
        \sum\limits_j s_j^{2(1-\alpha)}\abs{\braket{u_i}{v_j}}^2<\infty, \quad \forall i \text{ such that } r_i\neq 0
    \end{equation}
    \item \label{item:2-general}$\sum\limits_{i,j} r_i^{\alpha}s_j^{1-\alpha}\abs{\braket{u_i}{v_j}}^2 <\infty$.
\end{enumerate}
Furthermore, \begin{align}
    \label{eq:alpha-relative-entropy}
    D_{\alpha} (\rho \lvert\lvert \sigma) = \frac{1}{\alpha-1}\log \sum\limits_{i,j} r_i^{\alpha}s_j^{1-\alpha}\abs{\braket{u_i}{v_j}}^2,
\end{align}
whenever the Condition \ref{item:1-general} above is satisfied.
\end{thm}
\begin{proof}
By Lemma \ref{sec:gaussian-states-lem} and Theorem \ref{sec:gaussian-states-prop}, the operator $\rho^{\alpha/2}\sigma^{1-\alpha}\rho^{\alpha/2}$ is trace class if and only if the following two conditions are satisfied: 
\begin{enumerate}[label=(\roman*)]
    \item \label{item:1-general-proof} $u_i\in D(\sigma^{(1-\alpha)})$, $\forall i $ such that $r_i\neq 0$;
    \item \label{item:2-general-proof} $\int_0^\infty x\mu^{\rho^\alpha,\sigma^{(1-\alpha)}}(\dd x)<\infty$,
\end{enumerate}
where \ref{item:1-general-proof} ensures that the operator $\rho^{\alpha/2}\sigma^{1-\alpha}\rho^{\alpha/2}$ is well defined and \ref{item:2-general-proof} provides the value of a sum determining the trace. We will show that Condition \ref{item:1-general-proof} above is equivalent to (\ref{eq:condition-1})   and Condition \ref{item:2-general-proof} above is equivalent to  Condition \ref{item:2-general} in the statement of the present theorem. To this end, let $\xi^{\sigma^{1-\alpha}}$ denote the spectral measure corresponding to the operator $\sigma^{1-\alpha}$. Then \begin{align}
    \label{eq:spectral-measure-sigma}
\xi^{\sigma^{1-\alpha}}\{s_j^{1-\alpha}\} = \sum\limits_{\{r\in \BZ|s_j=s_r\}}\ketbra{v_r}.
\end{align} Now by spectral theory, for $i$ such that $r_i\neq 0$ in Condition \ref{item:1-general-proof} \begin{align*}
   &  u_i\in D(\sigma^{(1-\alpha)})\\
   & \Leftrightarrow \int\limits_0^\infty \abs{x}^2\mel{u_i}{\xi^{\sigma^{1-\alpha}}(\dd x)}{u_i}<\infty\quad \\
   & \Leftrightarrow \sum\limits_{j} \abs{s_j}^{2(1-\alpha)}\mel{u_i}{\xi^{\sigma^{1-\alpha}}\{s_j\}}{u_i}<\infty \\
 & \Leftrightarrow \sum\limits_{j} \abs{s_j}^{2(1-\alpha)}\mel{u_i}{\left(\sum\limits_{\{r\in \BZ|s_j=s_r\}}\ketbra{v_r}\right)}{u_i}<\infty\\
 & \Leftrightarrow \sum\limits_{j}\sum\limits_{\{r\in \BZ|s_j=s_r\}} \abs{s_j}^{2(1-\alpha)}\mel{u_i}{\left(\ketbra{v_r}\right)}{u_i}<\infty\\
 & \Leftrightarrow \sum\limits_{j}\sum\limits_{\{r\in \BZ|s_j=s_r\}} \abs{s_j}^{2(1-\alpha)}\abs{\braket{u_i}{v_r}}^2<\infty\\
 & \Leftrightarrow \sum\limits_{j} \abs{s_j}^{2(1-\alpha)}\abs{\braket{u_i}{v_j}}^2<\infty
\end{align*}
which proves (\ref{eq:condition-1}).
Now a similar computation as above using (\ref{item:7}), (\ref{eq:9}) and the fact that $\mu^{\rho^\alpha,\sigma^{(1-\alpha)}}$ is supported on the spectrum of $\sigma^{1-\alpha}$ shows that \begin{align*}
  \tr \rho^{\alpha/2}\sigma^{1-\alpha}\rho^{\alpha/2} &= \int_0^\infty x\mu^{\rho^\alpha,\sigma^{(1-\alpha)}}(\dd x)\\
  &= \sum\limits_j s_j^{1-\alpha}\sum\limits_ir_i^{\alpha}\mel{u_i}{\xi^{\sigma^{1-\alpha}}\{s_j^{1-\alpha}\}}{u_i}\\
  & = \sum\limits_{i,j}r_i^{\alpha}s_j^{1-\alpha}\abs{\braket{u_i}{v_j}}^2.
\end{align*}
Thus  Condition \ref{item:2-general-proof}
 is same as Condition \ref{item:2-general} and also (\ref{eq:alpha-relative-entropy}) is proved.
 \end{proof}
 \begin{rmks}\begin{enumerate}\item It may be noted that the Condition \ref{item:1-general} merely states that the operator  $\rho^{\alpha/2}\sigma^{1-\alpha}\rho^{\alpha/2}$ is well defined. Hence the Condition \ref{item:2-general} is the only condition necessary and sufficient for the finiteness of the Petz-R\'enyi $\alpha$-relative entropy between two states $\rho$ and $\sigma$, if we already know that the operator  $\rho^{\alpha/2}\sigma^{1-\alpha}\rho^{\alpha/2}$ is well defined.
     \item 
 In our definition of Petz-R\'enyi relative entropy between $\rho$
and $\sigma$ the finiteness of $D_\alpha(\rho||\sigma)$ implies (by Lemma \ref{sec:gaussian-states-lem}) that the operator $\rho^{\alpha/2}\sigma^{(1-\alpha)}\rho^{\alpha/2}$ is a bounded operator. In Chapter 2 of  \emph{Mathematical Foundations of Quantum Mechanics} 
\cite{vN2018}, John von Neumann defines the trace of the product of two unbounded positive operators $A,B$,  which we denote as $\operatorname{vN-Tr} \{AB\}$,  by \begin{align}\label{eq:vN-trace}
    \operatorname{vN-Tr} \{AB\} = \sum\limits_j \norm{\sqrt{A}\sqrt{B}u_j}^2,
\end{align}
where $\{u_j\}$ is an orthonormal basis of the Hilbert space contained
in the domain of $B$. Then the sum on the right side of (\ref{eq:vN-trace}) is independent of the choice of basis. We can define  \begin{align*}
    \operatorname{vN-D}_{\alpha}(\rho||\sigma) := \frac{1}{\alpha-1}\log \left(\operatorname{vN-Tr} \{\rho^{\alpha}\sigma^{(1-\alpha)}\}\right).
\end{align*}
 Now we compute the $\operatorname{vN-Tr} \{\rho^{\alpha}\sigma^{(1-\alpha)}\}$ using an orthonormal basis consisting of eigenvectors of $\sigma$ as in Theorem \ref{thm:the-formula-general},
  \begin{align*}
     \operatorname{vN-Tr} \{\rho^{\alpha}\sigma^{(1-\alpha)}\}  &=\sum\limits_{j} \norm{\rho^{\alpha/2}\sigma^{(1-\alpha)/2}v_j}^2\\
     & = \sum\limits_{j}
       s_j^{1-\alpha}\norm{\rho^{\alpha/2}v_j}^2\numberthis
       \label{eq:alternate-formula}\\
     & = \sum\limits_{j}
       s_j^{1-\alpha}\mel{v_j}{\rho^{\alpha}}{v_j}\\
     & = \sum\limits_{j}
       s_j^{1-\alpha}\mel{v_j}{\left(\sum\limits_ir_i^{\alpha}\ketbra{u_i}\right)}{v_j}\\
   & = \sum\limits_{i,j}r_i^{\alpha}s_j^{1-\alpha}\abs{\braket{u_i}{v_j}}^2.
   \end{align*}
    Thus the
   formula (\ref{eq:alpha-relative-entropy})  for the relative
   entropy remains the same in this case also. Since $\rho^{\alpha}\sigma^{(1-\alpha)}$ is a densely defined operator, we do not require Condition~1 of Theorem \ref{thm:the-formula-general} to hold in this case. The definition of trace according to von Neumann has the property, $\operatorname{vN-Tr} AB = \operatorname{vN-Tr} BA$, whenever $AB$ and $BA$ are well-defined. It is not clear to us (from the book \cite{vN2018}) whether von Neumann wants to define the trace only when  $AB$ and $BA$  both are well defined or not. Nevertheless, we have $ \operatorname{vN-Tr} \rho^{\alpha}\sigma^{1-\alpha} = \operatorname{vN-Tr} \sigma^{1-\alpha}\rho^{\alpha}$ whenever the latter  is well-defined, but in this case Condition~1  in Theorem \ref{thm:the-formula-general} is needed. For example, if $\alpha>1$ and $\sigma$ is any state, then 
   $\sigma^{1-\alpha}$ is an unbounded operator. Let $u\in \CK$ be such that $u\notin D(\sigma^{1-\alpha})$ and choose $\rho = \ketbra{u}$. In this case $\sigma^{1-\alpha}\rho^\alpha$ is not well defined.
   \end{enumerate}
 \end{rmks}
 
Now we prove one of the major results in this article. The reader is referred to Appendix \ref{sec:classical-renyi-devergence} for the basic definitions and properties related to the classical Petz-R\'enyi divergence.
\begin{thm}[Petz-R\'enyi relative entropy reduces to R\'enyi relative divergence]\label{thm:quantum-reduces-to-classical}
Let $\CK$ be a complex Hilbert space with  $\dim \CK = \abs{\mathcal{I}}$, where $\mathcal{I} =\{1,2,\dots,n\}$ or $\mathcal{I}=\BN$. 
Let $\rho$ and $\sigma$ be as in \eqref{eq:spectral-rho-sigma}.  
 Define $P$ and $Q$ on $\mathcal{I}\times \mathcal{I}$ such that,
\begin{align}
    \label{eq:P-and-Q}\begin{split}
  P(i,j)=  p_{ij} &= r_i\abs{\braket{u_i}{v_j}}^2,\\
  Q(i,j)=  q_{ij}& = s_j\abs{\braket{u_i}{v_j}}^2, \quad \forall (i,j)\in \mathcal{I}\times\mathcal{I}.
     \end{split}
\end{align} Then $P$ and $Q$ are probability measures on $\mathcal{I}\times \mathcal{I}$ such that \begin{equation}\label{eq:main-result}
  D_{\alpha}(\rho||\sigma) = D_{\alpha}(P||Q) ,\quad \forall \alpha\in (0,1)\cup(1,\infty),
\end{equation}
where the $D_\alpha$ on the left side of the equation above denotes the Petz-R\'enyi relative entropy of quantum states $\rho$ and  $\sigma$, and that on the right side is the classical R\'enyi relative divergence of probability distributions $P$ and $Q$.
\end{thm}
\begin{proof}
Since $\{u_i\}_{i}$ and $\{v_j\}_{j}$ are orthonormal bases, \[\sum_{i,j}p_{ij} = \sum_{i}r_i\sum_j\abs{\braket{u_i}{v_j}}^2 = \sum_{i}r_i = 1.\]
Similarly, $\sum_{i,j}q_{ij} = 1$. Hence $P$ and $Q$ are probability distributions on $\mathcal{I}\times \mathcal{I}$. Let $\mu$ be the counting measure on $\mathcal{I}\times \mathcal{I}$, then clearly $P\ll\mu$, $Q\ll\mu$, \begin{align*}
    \dv{P}{\mu}\,(i,j)=p_{ij}, \quad \textnormal{and} \quad \dv{Q}{\mu}\,(i,j)=q_{ij},
\end{align*} where $\dv{P}{\mu}$ and $\dv{Q}{\mu}$ are the respective Radon-Nikodym derivatives. Now by Definition \ref{defn:renyi-divergence-Kullback},  and equation (\ref{eq:alpha-relative-entropy}) in Theorem \ref{thm:the-formula-general},  for $\alpha\in (0,1)\cup (1,\infty)$ the R\'enyi divergence of $P$ from $Q$ is given by 
\begin{align*}
 D_{\alpha}(P||Q) &= \frac{1}{\alpha-1} \log \int_{\mathcal{I}\times \mathcal{I}} \left(\dv{P}{\mu}\right)^\alpha \left(\dv{Q}{\mu}\right)^{1-\alpha} \dd \mu   \\
&= \frac{1}{\alpha-1} \log \sum\limits_{i,j} \left(r_i\abs{\braket{u_i}{v_j}}^2\right)^\alpha\left(s_j\abs{\braket{u_i}{v_j}}^2\right)^{1-\alpha}\\
&=\frac{1}{\alpha-1} \log \sum\limits_{i,j} r_i^\alpha s_j^{1-\alpha}\abs{\braket{u_i}{v_j}}^2\\
&= D_{\alpha}(\rho||\sigma),
\end{align*}
where for $\alpha>1$, we adopt the conventions $0^{1-\alpha} = \infty$ and $0\cdot \infty = 0$ as in Definition \ref{defn:renyi-divergence-Kullback}
\end{proof}
\begin{rmk}\label{rmk:quamtum-classical-measurement}
In the setting of Theorem \ref{thm:quantum-reduces-to-classical} and its proof, for $i,j\in \mathcal{I}$ define $A_{ij} = \braket{u_i}{v_j}\ketbra{u_i}{v_j}$. Then $\sum_{i,j}A_{ij}A_{ij}^\dagger = I = \sum_{i,j}A_{ij}^\dagger A_{ij}$. Hence both $\{A_{ij}A_{ij}^\dagger\}$ and $\{A_{ij}^\dagger A_{ij}\}$ are POVM's. Furthermore, it may be noted that $p_{ij} = \tr \rho A_{ij}A_{ij}^\dagger$ and $q_{ij} =  \tr \sigma A_{ij}^\dagger A_{ij}$. Thus the probability measures $P$ and $Q$ in the previous theorem are precisely those measures that are obtained by measuring $\rho$ and $\sigma$, respectively in $\{A_{ij}A_{ij}^\dagger\}$ and  $\{A_{ij}^\dagger A_{ij}\}$.
\end{rmk}
Now we prove a lemma and two propositions which describe some relationships between the pairs $(\rho, \sigma)$ and $(P,Q)$  described in (\ref{eq:spectral-rho-sigma}) and (\ref{eq:P-and-Q}).
\begin{lem}\label{lem:support-condition}
Let $\rho$ and $\sigma$ be as in (\ref{eq:spectral-rho-sigma}). Then  $\supp \rho \subseteq \supp \sigma$ if and only if $s_j=0$ for some $j$ implies that for every $i$ at least one of the two quantities $\{\braket{u_i}{v_j}, r_i \}$ is equal to zero.  
\end{lem}
\begin{proof} ($\Rightarrow$) Assume $\supp \rho \subseteq \supp \sigma$.
The assumption that $s_j=0$ implies that $v_j\in \ker \sigma$. Now the condition $\braket{u_i}{v_j}\neq 0$ implies that $u_i\notin (\ker{\sigma})^{\perp} =\supp \sigma$. Since $\supp \sigma\supseteq\supp \rho$ and $u_i$ is an eigenvector of $\rho$ we see that \[u_i\notin\supp\rho \Rightarrow u_i\in \ker \rho \Rightarrow r_i=0.\]

($\Leftarrow$) We will prove that $\ker \sigma \subseteq \ker \sigma$. It is enough to prove that every $v_j$ for which $s_j = 0$ belong to $\ker \rho$. Fix such  $j$ such that $s_j=0$, by our assumption, $\braket{u_i}{v_j}=0$ for all $i$ such that $r_i\neq 0$ and hence $\sum_{\{i:r_i\neq 0\}}\braket{u_i}{v_j}\ket{u_i} = 0$. This means that the projection of $v_j$ to  $\supp \rho$ is the vector is the zero vector which means that $v_j\in \ker \rho$.
\end{proof}
\begin{prop}\label{prop:rho-equals-sigma-P-equals-Q}
Let $\rho$ and $\sigma$ be as in (\ref{eq:spectral-rho-sigma})  and let $P$ and $Q$ be as in (\ref{eq:P-and-Q}) then 
\[P=Q \Leftrightarrow \rho=\sigma.\]
\end{prop}
\begin{proof}Clearly, if $\rho=\sigma$ then $P=Q$. Now assume $P=Q$.
By definition \begin{align}\label{eq:P-equals-Q-implies}
    P=Q \Rightarrow r_i\abs{\braket{u_i}{v_j}}^2 = s_j\abs{\braket{u_i}{v_j}}^2,\quad \forall i,j.
\end{align}
Therefore, $s_j=0$ for some $j$ implies that for every $i$ at least one of the two quantities $\{\braket{u_i}{v_j}, r_i \}$ is equal to zero. Now by Lemma \ref{lem:support-condition}, $\supp\rho \subseteq \supp \sigma$. Also $\supp \sigma \subseteq \supp \rho$ by symmetry of the situation. Hence we have \[\supp\rho = \supp \sigma.\]
Therefore, $\ker\rho = \ker\sigma$. To prove that $\rho = \sigma$, we will prove that the nonzero eigenvalues and corresponding eigenspaces of $\rho$ and $\sigma$ are same.
Now fix $i_0\in \mathcal{I}$ such that $r_{i_0}\neq 0$. If $r_{i_0}\neq s_j$ for all $j$, the second equality in (\ref{eq:P-equals-Q-implies}) shows that $\braket{u_{i_0}}{v_j} =0$ for all $j$ such that $s_j\neq 0$. But this is impossible because $\{v_j|s_j\neq0\}$ is an orthonormal basis for $\supp \sigma$ and $0\neq u_{i_0}\in \supp \sigma$. Therefore, for each $i_0\in \mathcal{I}$ such that $r_{i_0}\neq 0$, there exists $j_0\in \mathcal{I}$ such that $r_{i_0}=s_{j_0}.$ Hence $\operatorname{sp}(\rho) \subseteq\operatorname{sp}(\sigma)$. A similar argument shows that $\operatorname{sp}(\sigma) \subseteq \operatorname{sp}(\rho)$. Thus we have  \[\operatorname{sp}(\rho) =\operatorname{sp}(\sigma).\]
Fix ${i_0}$ and ${j_0}$ such that $r_{i_0}=s_{j_0}$.
Let $R_{i_0}$ denote the eigenspace of $\rho$ corresponding to the eigenvalue $r_{i_0}$ and $S_{j_0}$ denote the eigenspace of $\sigma$ corresponding to the eigenvalue $s_{j_0}$. We will show that $R_{i_0}=S_{j_0}$, which will complete the proof since $r_{i_{0}}$ is an arbitrary non zero eigenvalue. It is enough to show the following two claims: \begin{enumerate}
    \item \label{claim:1} If $r_i=r_{i_0}$ for some $i$ then $u_i\perp v_j$ for all $j\in \mathcal{I}$ with $s_j\neq s_{j_0}$;
    \item If $s_j=s_{j_0}$ for some $j$ then $v_j\perp u_i$ for all $i\in \mathcal{I}$ with $r_i\neq r_{i_0}$.
\end{enumerate}
Proof of both the claims above are similar so we will prove \ref{claim:1} only. Fix $i$ such that $r_i=r_{i_0}$ and $j$ such that $s_j\neq s_{j_0}$. Since $r_{i_0}=s_{j_0}$, we have $r_i\neq s_j$. Thus by our assumption $\braket{u_i}{v_j}=0$, which proves \ref{claim:1}.
\end{proof}
 \begin{prop}\label{prop:support-condition-iff-absolute-continuity}
Let $\rho$ and $\sigma$ be as in (\ref{eq:spectral-rho-sigma})  and let $P$ and $Q$ be as in (\ref{eq:P-and-Q}), then \[\supp \rho \subseteq \supp \sigma \Leftrightarrow P\ll Q.\]
\end{prop}
\begin{proof}Assume $ P\nll Q$, we have
\begin{align*}
    P\nll Q &\Leftrightarrow \exists (i,j) \in \mathcal{I}\times \mathcal{I} \textnormal{ such that } \braket{u_i}{v_j} \neq 0, s_j=0,r_i\neq 0\\
    &\Leftrightarrow \exists (i,j) \in \mathcal{I}\times \mathcal{I} \textnormal{ such that } \braket{u_i}{v_j} \neq 0, v_j\in \ker \sigma, u_i \in \supp \rho\\
    &\Leftrightarrow \exists i\in \mathcal{I} \textnormal{ such that } u_i \in \supp \rho,  u_i \notin (\ker \sigma )^\perp\\
    &\Leftrightarrow \exists i\in \mathcal{I} \textnormal{ such that } u_i \in \supp \rho,  u_i \notin \supp \sigma\\
     &\Leftrightarrow \supp \rho \nsubseteq \supp \sigma
\end{align*}
\end{proof}
\subsubsection{The Limiting Cases}
The definition of Petz-R\'enyi $\alpha$-relative entropy excludes the values $0$, $1$ and $\infty$ of $\alpha$. Nevertheless, we can give meaning to the  entropic quantities corresponding to these values of $\alpha$ and they are important in applications too \cite{Datta2009, muller-tomamichel-etal-2013, Datta_Leditzky_2014}. In this section, we study the limits at $0$ and $\infty$ of $D_{\alpha}(\cdot||\cdot)$. The case $\alpha= 1$ is requires more analysis and we discuss that in detail in Section \ref{sec:vN-entropy}.
First we prove that the  Petz-R\'enyi $\alpha$-relative entropy is nondecreasing in $\alpha$, which will help us to extend the definition of $D_{\alpha}$ to the values $0, 1$ and $\infty$.
\begin{thm} \label{thm:3-veh-quantum}
For $\alpha\in (0,1)\cup(1,\infty)$ the  Petz-R\'enyi entropy, $D_{\alpha}(\rho||\sigma)$ is nondecreasing in $\alpha$.
\end{thm}
\begin{proof}
This is an easy consequence of Theorem \ref{thm:quantum-reduces-to-classical} and Theorem \ref{thm:3-veh}.
\end{proof}

 Theorem \ref{thm:3-veh-quantum} enables us to extend the definition of $D_{\alpha}(\rho||\sigma)$ to the values $\alpha = 0,1$ and $\infty$ as in the following definition.
\begin{defn}\label{defn:extended-orders-quantum}
 The Petz-R\'enyi relative entropies of orders $0, 1$ and $\infty$ are defined as \begin{align*}
     D_0(\rho||\sigma) &= \lim\limits_{\alpha \downarrow 0} D_{\alpha}(\rho||\sigma),\\
     D_1(\rho||\sigma) &= \lim\limits_{\alpha \uparrow 1} D_{\alpha}(\rho||\sigma),\\
     D_{\infty}(\rho||\sigma) &= \lim\limits_{\alpha \uparrow \infty} D_{\alpha}(\rho||\sigma).
 \end{align*}
\end{defn}
With the definition above, we have the following corollary.
\begin{cor}
For $\alpha\in [0,\infty]$, the function $\alpha\mapsto D_{\alpha}(\rho||\sigma)$ is nondecreasing and thus \[D_0(\rho||\sigma)\leq D_{\alpha}(\rho||\sigma)\leq D_{\infty}(\rho||\sigma), \quad \forall \alpha\geq0.\]
\end{cor}
\begin{rmk} Following the notations used in \cite{Datta2009},
the quantity $D_0(\rho||\sigma)$  may also be written as $D_{\operatorname{min}}(\rho||\sigma)$.
\end{rmk}
If $\rho = \sum_ir_i\ketbra{u_i}{u_i}$ is a spectral decomposition of $\rho$, where $\{u_i\}$ is an orthonormal basis of $\CK$, then
\begin{equation}\label{eq:support-formula}
 \supp\rho = \overline{\spn}\{u_i|r_i\neq 0 \}.   
\end{equation} 
In \cite{Datta2009}, Datta observes in the finite dimensional setting that,  \begin{align}\label{eq:datta-leditzky}
    D_0(\rho||\sigma) = -\log \tr \Pi_{\rho}\sigma,
\end{align}
where $\Pi_{\rho}$ is the projection onto the support of $\rho$, i.e., by keeping the notations of Theorem \ref{thm:the-formula-general}, \[\Pi_{\rho} = \textnormal{Projection onto the }\spn\{u_i\,|\,r_i\neq 0\}.\] In the finite dimensions, this result  follows from our  formula (\ref{eq:alpha-relative-entropy}) as well, because  we have \begin{align}\label{eq:datta-leditzky-computation}
    \begin{split}
       D_0(\rho||\sigma)&= \lim_{\alpha\rightarrow0}\frac{1}{\alpha-1}\log \sum\limits_{i,j} r_i^{\alpha}s_j^{1-\alpha}\abs{\braket{u_i}{v_j}}^2\\
       & = -\log \sum\limits_{\{i,j|r_i\neq 0\}}s_j^{}\abs{\braket{u_i}{v_j}}^2\\
       &=-\log \tr \sum_{\{i|r_i\neq 0\}}\ketbra{u_i}{u_i}\sum_{j}s_j\ketbra{v_j}{v_j}\\
       &= -\log \tr \Pi_{\rho}\sigma.
    \end{split}
\end{align}
 Now we prove (\ref{eq:datta-leditzky}) in the infinite dimensional situation. A priori, the computation in (\ref{eq:datta-leditzky-computation}) cannot go through in  infinite dimensions because we need a limit theorem to pass the limit through the infinite sum. Nevertheless, Theorem \ref{thm:4-veh} helps us to prove the desired result and the following proof works in both  finite and infinite dimensional setting.
 \begin{thm}\label{thm:4-veh-quantum}
 The Petz-R\'enyi relative entropy satisfies, \begin{align}\label{eq:datta-leditzky-infinite}
    D_0(\rho||\sigma) = -\log \tr \Pi_{\rho}\sigma,
\end{align}
where $\Pi_{\rho}$ is the projection onto  $\supp \rho$.
 \end{thm}
 \begin{proof}
 Keeping the notations of Theorem \ref{thm:the-formula-general}, we have by Theorems \ref{thm:quantum-reduces-to-classical} and \ref{thm:4-veh},
 \begin{align*}
      D_0(\rho||\sigma) &= -\log (Q ({\{p_{ij}>0\}}))\\
      & = \sum_{\{i,j|r_i>0, \braket{u_i}{v_j}\neq0\}}s_j\abs{\braket{u_i}{v_j}}^2\\
      & = -\log \sum\limits_{\{i,j|r_i\neq 0\}}s_j^{}\abs{\braket{u_i}{v_j}}^2\\
       &=-\log \tr \sum_{\{i|r_i\neq 0\}}\ketbra{u_i}{u_i}\sum_{j}s_j\ketbra{v_j}{v_j}\\
       &= -\log \tr \Pi_{\rho}\sigma.
 \end{align*}
 \end{proof}
 \begin{rmk}
 On a related note, it may be recalled that a sandwiched R\'enyi relative entropy, $\tilde{D}_{\alpha}$ was introduced by M\"{u}ller-Lennert et al. in \cite{muller-tomamichel-etal-2013}. Furthermore, Datta and Leditzky in their Theorem 1 of \cite{Datta_Leditzky_2014} proved that \[\lim_{\alpha\rightarrow 0 }\tilde{D}_{\alpha} = D_0(\rho||\sigma),\]
 whenever $\supp \rho =\supp \sigma$.
 \end{rmk}
 \begin{thm}\label{thm:6-veh-quantum}
Let $\rho$ and $\sigma$ be as in (\ref{eq:spectral-rho-sigma}). Then 
\begin{equation}
 \label{eq:limit-infty-quantum}
 D_{\infty}(\rho||\sigma) = \log \underset{}{\sup}\left\{\frac{r_i}{s_j}\,:\,\braket{u_i}{v_j}\neq 0\right\},
\end{equation}
with the conventions that $0/0 = 0$ and $x/0 = \infty$ if $x>0$. 
 \end{thm}
 \begin{proof} Let $P$ and $Q$ be as in (\ref{eq:P-and-Q}).
By the definition  of $D_{\infty}(\rho||\sigma)$ and Theorem \ref{thm:quantum-reduces-to-classical}, we have 
\begin{align*}
 D_{\infty}(\rho||\sigma)&= \lim_{\alpha\rightarrow \infty}D_{\alpha}(\rho||\sigma)= \lim_{\alpha\rightarrow \infty}D_{\alpha}(P||Q) = D_{\infty}(P||Q)
\end{align*}
But by equation (\ref{eq:limit-infty-countable-sample-space}), \begin{align*}
  D_{\infty}(P||Q) &= \log \underset{(i,j)\in \mathcal{I}\times \mathcal{I}}{\sup}\,\frac{r_i \abs{\braket{u_i}{v_j}^2}}{s_j\abs{\braket{u_i}{v_j}^2}}\\
  &=\log \underset{}{\sup}\left\{\frac{r_i}{s_j}\,:\,\braket{u_i}{v_j}\neq 0\right\}
\end{align*}
with the conventions that $0/0 = 0$ and $x/0 = \infty$ if $x>0$.
 \end{proof}
 In the following Corollary, we remove the condition $\braket{u_i}{v_j}\neq 0$ from the formula for $D_{\infty}(\rho||\sigma)$ obtained in the previous theorem. But this will result in imposing an extra condition on the $\rho$ and $\sigma$.
 \begin{cor}\label{cor:d-infty}
 Let $\rho$ and $\sigma$ be as in (\ref{eq:spectral-rho-sigma}). Assume further that \begin{align}\label{eq:D-infty-hypothesis}
     \braket{u_i}{v_j}=0 \text{ for some } (i,j) \Rightarrow r_i = 0,
 \end{align}
 then 
 \begin{align}
     D_{\infty}(\rho||\sigma)  =\begin{cases}
       \log \norm{\rho}\norm{\sigma^{-1}}, \quad &\text{ when } \dim \CK<\infty;\\
       \infty,\quad &\text{ when } \dim \CK=\infty,
       \end{cases}
 \end{align}
where $\norm{\cdot}$ denotes the operator norm,   $\norm{\sigma^{-1}}$ is defined to be infinity if $\sigma$ is not invertible, and $\log \infty:=\infty$.
 \end{cor}
 \begin{proof}
 Notice that in Equation (\ref{eq:limit-infty-quantum}) we can  remove the restriction
 $\braket{u_i}{v_j} \not = 0$, since when $\braket{u_i}{v_j} = 0$, then by the hypothesis (\ref{eq:D-infty-hypothesis}), the fraction $\frac{r_i}{s_j}$ is equal to $0$ and it does
 not contribute to the $\sup \{ \frac{r_i}{s_j}: (i,j) \in \mathcal{I}\times \mathcal{I}\}$. Now
 \begin{align*}
   D_{\infty}(\rho||\sigma)&= \log \underset{\underset{}{(i,j)\in \mathcal{I}\times \mathcal{I}}}{\sup}\,\frac{r_i }{s_j}.
   \end{align*}
   Note that if $\dim \CK = \infty$ then $s_j=0$ for some $j$ or there is a subsequence of $s_j$'s which converges to zero. In both cases we get $\log \underset{\underset{}{(i,j)\in \mathcal{I}\times \mathcal{I}}}{\sup}\,\frac{r_i }{s_j} = \infty$. On the other hand, if $\dim \CK <\infty$ we have 
   \begin{align*}
   D_{\infty}(\rho||\sigma)&= \log \underset{\underset{r_i\neq0}{(i,j)\in \mathcal{I}\times \mathcal{I}}}{\sup}\,\frac{r_i }{s_j}\\
   &=  \underset{\underset{r_i\neq0}{(i,j)\in \mathcal{I}\times \mathcal{I}}}{\sup} \log (\frac{r_i}{s_j})\\
   &=  \underset{\underset{r_i\neq0}{(i,j)\in \mathcal{I}\times \mathcal{I}}}{\sup} \left(\log r_i+\log{{s_j^{-1}}}\right) \\
   &= \underset{r_i\neq0}{\sup} \{\log r_i\} +\underset{}{\sup} \{\log {s_j^{-1}}\}\\
   &=\log \left(\underset{r_i\neq0}{\sup} \{r_i\}\right)+\log\left( \underset{}{\sup} \left\{{s_j^{-1}}\right\}\right)\\
   &= \log \norm{\rho} +\log \norm{\sigma^{-1}}\\
   &= \log \norm{\rho}\norm{\sigma^{-1}}.
 \end{align*} 
 \end{proof}
 \begin{rmk}\label{rmk:camparison-datta}Recall the definition of $D_{\operatorname{max}}(\rho||\sigma)$ from equation (7) in \cite{Datta2009},   \begin{align*}
     D_{\operatorname{max}}(\rho||\sigma):= \log \norm{\sigma^{-1/2}\rho\sigma^{-1/2}},
 \end{align*}
 whenever $\supp\rho \subseteq\supp \sigma$. When $\supp \rho \subseteq \supp \sigma$, the  equation (31) of \cite{Datta2009} notes that \begin{align}
     \label{eq:datta-d-max}
     D_{\operatorname{max}}(\rho||\sigma) \leq \log \sup\left\{\frac{1}{s_j}\,:\, s_j\neq 0\right\}.
 \end{align}
Let $\operatorname{sp}(A)$ denote the spectrum of an operator $A$.  Note that \[\norm{\sigma^{-1/2}\rho\sigma^{-1/2}}=\sup \operatorname{sp}(\sigma^{-1/2}\rho\sigma^{-1/2})=\sup \operatorname{sp}(\rho\sigma^{-1})\leq\norm{\rho\sigma^{-1}}\leq \norm{\rho} \norm{\sigma^{-1}}.\]
Therefore, when (\ref{eq:D-infty-hypothesis}) is satisfied and $\sigma$ is invertible, Corollary \ref{cor:d-infty} shows that \begin{align}\label{eq:comparison-datta}
     D_{\operatorname{max}}(\rho||\sigma):= \log \norm{\sigma^{-1/2}\rho\sigma^{-1/2}}&\leq \log \norm{\rho} \norm{\sigma^{-1}} = \log \norm{\rho}+\log \norm{\sigma^{-1}}\nonumber\\&\leq \log \norm{\sigma^{-1}}  = \log\sup\left\{\frac{1}{s_j}\,:\, s_j\neq 0\right\}.
\end{align} Thus equation (\ref{eq:comparison-datta}) above improves  (\ref{eq:datta-d-max}) under the extra assumption (\ref{eq:D-infty-hypothesis}) and $\sigma$ is invertible.
 \end{rmk}
 Next example shows that the expressions in Theorem \ref{thm:6-veh-quantum} and Corollary \ref{cor:d-infty} for $D_{\infty}(\rho||\sigma)$ are different in general.
 \begin{eg}
   Let $\dim \CK  = 2$ and let $\{u_1,u_2\}$ be any orthonormal basis of $\CK $. Let \begin{align*}
       \rho = \sigma =  \frac{2}{3}\ketbra{u_1}+\frac{1}{3}\ketbra{u_2}.
   \end{align*}  
   Then $r_1=s_1=2/3$, $r_2=s_2=1/3$ and \[\left\{\frac{r_i}{s_j}\,:\,\braket{u_i}{v_j}\neq 0\right\}  =\{1\}, \quad \left\{\frac{r_i}{s_j}\,:\, r_i\neq0\right\}= \{2,1\}. \]
 Hence \[\log \underset{\underset{r_i\neq0,}{(i,j)\in \mathcal{I}\times \mathcal{I}}}{\sup}\,\frac{r_i }{s_j}=\log 2=1 > 0 = \log \underset{}{\sup}\left\{\frac{r_i}{s_j}\,:\,\braket{u_i}{v_j}\neq 0\right\}.  \]
  \end{eg}
\subsubsection{von Neumann Relative Entropy} \label{sec:vN-entropy} In this section, we discuss the limiting case $D_1(\rho||\sigma)$  for the Petz-R\'enyi $\alpha$-relative entropy. Theorem \ref{thm:von-Neumann-entropy} shows that this limit is the von Neumann entropy.  The result is straight forward in the finite dimensional setting. But the infinite dimensional situation is a bit tricky because of reasons we shall see below. One definition of the von Neumann entropy in the infinite dimensional setting is  Araki's definition \cite{araki1976relative, araki1977relative} for general semifinite von Neumann Algebra, using the heavy machinery of relative modular operator. To avoid this, we will show that it makes sense to define the von Neumann entropy as $D_1(\rho||\sigma)$ itself when $\CK $ is infinite dimensional. We begin with the definition of von Neumann relative entropy in the finite dimensional setting.
 
 
 In the finite dimensional situation, the von Neumann relative entropy $D(\rho||\sigma)$ of two states $\rho$ and $\sigma$ are defined as \begin{align*}
     D(\rho||\sigma) = \begin{cases}
       \tr\rho(\log\rho-\log\sigma), &\textnormal{ when } \supp \rho \subseteq \supp \sigma;\\
       \infty, & \textnormal{otherwise.}
     \end{cases}
 \end{align*}
Let $\dim \CK  = n<\infty$, and let $\rho = \sum_{i=1}^{n}r_i \ketbra{u_i}, \sigma = \sum_{j=1}^{n}s_j \ketbra{v_j} $, where $\{u_i\}_{i=1}^{n}$, $\{v_j\}_{j=1}^{n}$ are orthonormal bases of $\CK $. The operators  $\log\rho$ and $\log \sigma$ have to be appropriately understood when $\rho$ and $\sigma$ have a  non trivial kernel. We may define
\begin{align}\begin{split}\label{eq:log-finite-dim}
    \log \rho &=  \sum_{r_i\neq0}\log r_i \ketbra{u_i} \textnormal{ on } \supp \rho\\
    \textnormal{and }\log \sigma &=  \sum_{s_j\neq0}\log s_j \ketbra{v_j} \textnormal{ on } \supp \sigma.
    \end{split}
\end{align} 
 Under the assumption $\supp \rho \subseteq \supp \sigma$, the operator $(\log \rho -\log \sigma)$ is well defined on $\supp \rho$. Now we describe another way to make sense of the definition of the von Neumann entropy. With the convention $0\cdot\infty=0$, the operator $\rho\log\rho$ can be defined as  \[\rho\log\rho =\sum_{j=1}^{n}r_i\log r_i \ketbra{u_i}. \]
Also, $\rho\log\sigma$ can be defined as \begin{align}\label{eq:rholog-sigma-finite-dim}
    \rho\log \sigma = \sum_{i,j=1}^n r_i\log s_j\braket{u_i}{v_j}\ketbra{u_i}{v_j},\quad\textnormal{whenever } \supp \rho \subseteq \supp \sigma,
\end{align}
because of Lemma \ref{lem:support-condition}.


We wish to set up a definition of $D(\rho||\sigma)$ which  works well both in the finite  and  infinite dimensional setting.  To this end, we first define the operator $-\log \sigma$ on a dense subspace of $\supp \sigma$ by the spectral decomposition
\begin{align}\label{eq:log-sigma}
    -\log \sigma = \sum\limits_{\{j\in \mathcal{I}\,|\, s_j\neq 0\}}-\log s_j\ketbra{v_j},
\end{align}
where $v_j$ is considered as an element of the Hilbert space $\supp \sigma$.
Note that, by the spectral theorem the domain of $-\log \sigma$ is given by \begin{align}\label{eq:domain-logsigma}
    D(-\log \sigma) = \{v\in \supp \sigma\,|\, \sum_{\{j\in \mathcal{I}\,|\,s_j\neq 0\}} \abs{\log s_j}^2\abs{\braket{v}{v_j}}^2<\infty\}\subseteq \supp \sigma.
\end{align}
    It is easy to see that $\spn\{v_j\,:\, s_j\neq 0\}\subseteq D(-\log \sigma)$.
Now $(-\log \sigma)$ is a positive (selfadjoint) operator on $D(-\log \sigma)\subseteq \supp\sigma$.  A similar definition applies to $\log \rho$ also, and $\log \rho$ defined on $D(\log \rho)\subseteq \supp \rho$ is a selfadjoint operator with spectrum on the negative half of the real line. It may be observed that $\log \rho $ (or $-\log \sigma)$ as defined above is bounded if and only if $\rho$ (or $\sigma$) has only a finite number of nonzero eigenvalues.
 \begin{eg}[$\supp \rho \subseteq \supp \sigma$, but the operator $(\log\rho -\log\sigma)$ does not make sense and $\rho^{1/2}(-\log\sigma)\rho^{1/2}$  is not densely defined]\label{eg:counter-example} \label{eg:support-condition-not-enough}
Let $\{v_j\}_{j=1}^{\infty}$ be an orthonormal basis of $\CK $. Take
\begin{align*}
    \sigma&= \sum_{j=1}^{\infty} 2^{-j}\ketbra{v_j}\\
    \rho &= \ketbra{u_1}, \textnormal{ where }  u_1 =\left(\sum_{j=1}^{\infty}\frac{1}{j^2}\right)^{-1/2}\sum_{j=1}^{\infty}\frac{1}{j}\ket{v_j}.
\end{align*}  Then $\supp \sigma = \CK $ and
 the condition $\supp \rho \subseteq \supp \sigma$ is trivially satisfied. Note that
 \[\sum_{j=1}^{\infty}\abs{\log 2^{-j}}^2\abs{\braket{u_1}{v_j}}^2 =\left(\sum_{j=1}^{\infty}\frac{1}{j^2}\right)^{-1/2} \sum_{j=1}^{\infty}\frac{j^2}{j^2}=\infty, \] so by (\ref{eq:domain-logsigma}),
 $u_1\notin D(-\log \sigma)$. But for any $ v\in \CK $ with $v\neq 0$,  $\rho^{1/2}v = \braket{u_1}{v}\ket{u_1}$ and hence $\rho^{1/2}v\in D(-\log \sigma)$ if and only if $v\in \ker \rho$. Therefore, \begin{enumerate}
     \item $ D(\rho^{1/2}(-\log\sigma) \rho^{1/2}) = \ker \rho$ and hence it is not densely defined;
     \item $\log \rho$ is defined only on $\spn \{u_1\}$ and $\log \rho =0$ on $\spn \{u_1\}$;
     \item $D(\log \rho -\log \sigma) =\{0\}$;
     \item $\rho^{1/2}(-\log\sigma) \rho^{1/2}=0$ on $\ker \rho$ and $\rho^{1/2}(-\log\rho)\rho^{1/2} =0 $.
 \end{enumerate}   
  Thus we have   \begin{align}
   \tr \rho^{1/2}(-\log\sigma)\rho^{1/2}-\tr\rho^{1/2}(-\log\rho)\rho^{1/2} = 0.  
 \end{align} 
 \end{eg}

  Let $\CK $ be a finite or infinite dimensional Hilbert space. Let  $\rho $ and  $\sigma$ be as in (\ref{eq:spectral-rho-sigma}).
 Assume that \begin{align}
    \label{eq:support-condition-infinite-dim}
    \spn \{u_i\,|\,r_i\neq 0 \}\subseteq \spn \{v_j\,|\,s_j\neq 0\}.
\end{align}
 Now by the first half of Theorem \ref{sec:gaussian-states-prop}, $\rho^{1/2}(-\log \sigma)\rho^{1/2}$ first defined on the dense domain $\spn\{u_i\}_{i\in \mathcal{I}} $ extends as a positive (selfadjoint) operator on a dense domain in $\CK $. Next theorem shows that our definition of the von-Neumann entropy coincides with the standard definition of the same in the finite dimensional setting, and also generalizes nicely to the infinite dimensional setting, whenever the quantities involved do make sense.
  \begin{thm}\label{thm:von-Neumann-entropy}Let $\rho$ and $\sigma$ be as in (\ref{eq:spectral-rho-sigma}).
Then 
 \begin{align}\label{eq:vn-entropy}
     \lim_{\alpha\uparrow 1} D_{\alpha}(\rho||\sigma)= \begin{cases}
     \tr \rho^{1/2}(-\log\sigma)\rho^{1/2}-\tr\rho^{1/2}(-\log\rho)\rho^{1/2},  &\text{when } (\ref{eq:support-condition-infinite-dim}) \text{ is satisfied} \\
      &\text{and at least one of the  }\\
   &\text{quantities }\\
   &\tr \rho^{1/2}(-\log\sigma)\rho^{1/2}\\
   &\textnormal{or } \\&\tr \rho^{1/2}(-\log\rho)\rho^{1/2} \text{is finite;}\\
   &\\
       \infty, & \text{when } \supp \rho \nsubseteq \supp \sigma.\\
     \end{cases}
 \end{align}
 In particular, if $\dim \CK <\infty $,  then \begin{align}
   \lim_{\alpha\uparrow 1} D_{\alpha}(\rho||\sigma) =\begin{cases}
       \tr \rho(\log \rho -\log \sigma),& \text{ when } \supp \rho \subseteq \supp \sigma;\\
       \infty, & \text{otherwise}.
     \end{cases} 
 \end{align}
 Moreover, if $D_1(\rho||\sigma) = \infty$ or there exists $\beta>1$ such that $D_{\beta}(\rho||\sigma)<\infty$, then also \begin{align}\label{eq:vn-entropy-equals-right-limit}
     D_1(\rho||\sigma)  = \lim\limits_{\alpha\downarrow 1}D_{\alpha}(\rho||\sigma).
 \end{align}
 \end{thm}
 \begin{proof}
Assume first that (\ref{eq:support-condition-infinite-dim}) is satisfied. By the remarks after (\ref{eq:support-condition-infinite-dim}) the operators $\rho^{1/2}(-\log\sigma)\rho^{1/2}$ and $\rho^{1/2}(-\log \rho)\rho^{1/2}$ are densely defined positive operators. We have
\begin{align*}
    \tr \rho^{1/2}(-\log\sigma)\rho^{1/2} &= \sum\limits_{i}\mel{u_i}{\rho^{1/2}(-\log\sigma)\rho^{1/2}}{u_i}\\
    &=\sum\limits_{i}r_i\mel{u_i}{(-\log\sigma)}{u_i}\\
     &=\sum\limits_{i}\sum\limits_{\{j\,|\, s_j\neq 0\}}r_i(-\log s_j)\mel{u_i}{(\ketbra{v_j})}{u_i}\\
      &=\sum\limits_{i}\sum\limits_{\{j\,|\, s_j\neq 0\}}-r_i(\log s_j)\abs{\braket{u_i}{v_j}}^2.
\end{align*}
By Remark \ref{rmk:support-condition}, Lemma \ref{lem:support-condition} is satisfied in our case. Now by using the convention $0\cdot\infty = 0$, we can sum over all possible $j$'s in the last sum above and write  \begin{align}\label{eq:first-term-vn-entropy}
    \tr \rho^{1/2}(-\log\sigma)\rho^{1/2}= \sum\limits_{i,j}-r_i(\log s_j)\abs{\braket{u_i}{v_j}}^2.
\end{align}
Similarly, we also have $\tr \rho^{1/2}(-\log\rho)\rho^{1/2}= \sum\limits_{i}-r_i(\log r_i)\abs{\braket{u_i}{v_j}}^2.$ Since $\sum_j\abs{\braket{u_i}{v_j}}^2 = 1$ for all $i$, we can write
\begin{align}\label{eq:second-term-vn-entropy}
    \tr \rho^{1/2}(-\log\rho)\rho^{1/2}= \sum\limits_{i,j}-r_i(\log r_i)\abs{\braket{u_i}{v_j}}^2.
\end{align}
By (\ref{eq:first-term-vn-entropy}) and (\ref{eq:second-term-vn-entropy}), we have  \begin{align}\label{eq:vn-entropy-formula-before-clubbing the terms}
 & \tr \rho^{1/2}(-\log\sigma)\rho^{1/2}-\tr\rho^{1/2}(-\log\rho)\rho^{1/2}\nonumber\\
 &\phantom{...........}= \sum\limits_{i,j}-r_i(\log s_j)\abs{\braket{u_i}{v_j}}^2 -\sum\limits_{i,j}-r_i(\log r_i)\abs{\braket{u_i}{v_j}}^2\nonumber\\
    &\phantom{...........}=\sum\limits_{i,j} r_i(\log r_i)\abs{\braket{u_i}{v_j}}^2-\sum\limits_{i,j}r_i(\log s_j)\abs{\braket{u_i}{v_j}}^2.\end{align}
    If at least one of  the quantities $\tr \rho^{1/2}(-\log\sigma)\rho^{1/2}$ or $\tr\rho^{1/2}(-\log\rho)\rho^{1/2}$ are  finite,we can combine the two summations above to write,\[\tr \rho^{1/2}(-\log\sigma)\rho^{1/2}-\tr\rho^{1/2}(-\log\rho)\rho^{1/2} = \sum\limits_{i,j} r_i(\log r_i)\abs{\braket{u_i}{v_j}}^2-r_i(\log s_j)\abs{\braket{u_i}{v_j}}^2.\]  By using Lemma \ref{lem:support-condition}, and the conventions $\frac{0}{0}=0$, $\frac{x}{0}=\infty$ if $x>0$, $0\cdot\infty = 0$,   to get  
    \begin{align*}
       \tr \rho^{1/2}(-\log\sigma)\rho^{1/2}-\tr\rho^{1/2}(-\log\rho)\rho^{1/2}
    & =\sum_{i,j}r_i\abs{\braket{u_i}{v_j}}^2\log\frac{r_i}{s_j}\\
    &= \sum_{i,j}r_i\abs{\braket{u_i}{v_j}}^2\log\frac{r_i\abs{\braket{u_i}{v_j}}^2}{s_j\abs{\braket{u_i}{v_j}}^2}\\
    &= D(P||Q),
\end{align*}where $D(P||Q)$ is the Kullback-Leibler divergence   of $P$ and $Q$,
where $P$ and $Q$ are as in (\ref{eq:P-and-Q}).  Now by 
Theorem \ref{thm:5-veh} we get
\begin{align*}
    D_1(\rho||\sigma) = D(P||Q) = \tr \rho^{1/2}(-\log\sigma)\rho^{1/2}-\tr\rho^{1/2}(-\log\rho)\rho^{1/2} 
\end{align*}
whenever (\ref{eq:support-condition-infinite-dim}) is satisfied and at least one of the quantities $\tr \rho^{1/2}(-\log\sigma)\rho^{1/2}$ or $\tr\rho^{1/2}(-\log\rho)\rho^{1/2}$ is finite.


To prove the second equality in (\ref{eq:vn-entropy}), assume that $\supp \rho \nsubseteq \supp \sigma$. Then Proposition \ref{prop:support-condition-iff-absolute-continuity} shows that $P\nll Q$, and hence  by equation (\ref{eq:kullback-leibler-equals-infty}), and Theorem \ref{thm:5-veh} the Kullback-Leibler divergence $D(P||Q) = \infty = \lim_{\alpha\uparrow} D_\alpha(P||Q) = D_1(\rho||\sigma)$ in this case.

Finally, equation (\ref{eq:vn-entropy-equals-right-limit}) follows from Theorem \ref{thm:5-veh} and Remark \ref{rmk:5-veh-quantum}
\end{proof}
 \begin{rmk}
\label{rmk:support-condition}
 Note that (\ref{eq:support-condition-infinite-dim}) is same as the condition $\supp\rho\subseteq\supp\sigma$ when $\dim \CK $ is finite. Also, the condition (\ref{eq:support-condition-infinite-dim}) implies $\supp\rho\subseteq\supp\sigma$. Hence Lemma \ref{lem:support-condition} holds  in this case as well.
 \end{rmk}
\begin{eg}\label{eg:infty-infty} If $\rho$ is any state such that $\tr \rho^{1/2}(-\log\rho)\rho^{1/2} = -\sum_{\{i\,:\, r_i\neq 0\}} r_i \log r_i= \infty$, then $ \tr \rho^{1/2}(-\log\rho)\rho^{1/2}-\tr\rho^{1/2}(-\log\rho)\rho^{1/2}= \infty -\infty$. A specific example of such a state is obtained by taking $\rho = \sum_{i}r_i\ketbra{u_i}{u_i}$ where $\{u_i\}$ is an orthonormal basis and  $r_i=\left(\sum_j \frac{1}{j(\log j)^2} \right)^{-1} \frac{1}{i(\log i)^2}$.
\end{eg}
Example \ref{eg:infty-infty} suggests that it is impossible to define the von-Neumann entropy in the infinite dimensional setting  by   $D(\rho||\sigma) = \tr \rho^{1/2}(-\log\sigma)\rho^{1/2}-\tr\rho^{1/2}(-\log\rho)\rho^{1/2}$, even when $(\ref{eq:support-condition-infinite-dim})$ is satisfied because this can lead to $D(\rho||\sigma)= \infty-\infty$,  in certain situations.  
Finally, in the light of Theorem \ref{thm:von-Neumann-entropy} we propose the following definition for von Neumann relative entropy.
 \begin{defn}\label{defn:vn-entropy}
 Define the von-Neumann relative entropy $D(\rho||\sigma)$ of two states $\rho$ and $\sigma$ by
 \begin{align}
     D(\rho||\sigma) = D_1(\rho||\sigma):= \lim_{\alpha \uparrow 1}D_{\alpha}(\rho||\sigma).
 \end{align}
 \end{defn}
 \begin{rmk}
 \label{rmk:5-veh-quantum} Let $\rho$ and $\sigma$ be as in (\ref{eq:spectral-rho-sigma})  and let $P$ and $Q$ be as in (\ref{eq:P-and-Q}).   By Theorem \ref{thm:5-veh} and Theorem \ref{thm:quantum-reduces-to-classical}, the von-Neumann entropy of $\rho$ and $\sigma$ is  the Kullback-Leibler divergence of $P$ and $Q$, i.e., we have \begin{align}\label{eq:5-veh-quantum} 
     D(\rho||\sigma) = D(P||Q).
 \end{align}
 \end{rmk} 
\subsubsection{Example: Araki's Relative Entropy is different from Information Theoretic Relative Entropy when \texorpdfstring{\ensuremath{\dim \CK  =\infty}}{}}\label{sec:counter-eg}
In this section we discuss a counter example to  Theorem 20 of \cite{luczak2019relative} thus disproving that result.
Consider $\rho$ and $\sigma$ as in  Example \ref{eg:counter-example}. It was shown that $\supp \rho \subseteq \supp \sigma$ and \[\tr \rho^{1/2}(\log \rho) \rho^{1/2}-\tr\rho^{1/2}(\log \sigma)\rho^{1/2}=0.\] We take the semifinite von Neumann algebra  $\B{\CK }$. If $\eta\in \B{\CK }$ is a density operator i.e., $\eta\geq 0$, and $\tr \eta =1$, then $\eta$ can be considered as a normal state $\varphi_{\eta}$ on the von Neumann algebra $\B{\CK }$ via its action \[X\overset{\varphi_{\eta}}{\longmapsto} \tr \eta X, \quad \forall X\in \B{\CK }.\] We will show that Araki's relative entropy $S(\varphi_{\rho}||\varphi_\sigma)$, of $\rho$ and $\sigma$ as in Example \ref{eg:counter-example} is\[ S(\varphi_{\rho}||\varphi_\sigma) = \infty\] 
in this case. 

The definition of Araki's relative entropy uses the notion of relative modular operator. A quick introduction to the description of relative modular operator in our setting is provided in Appendix \ref{appendix:relative-modular}. Let $\Bt{\CK }$ be the Hilbert space of Hilbert-Schmidt operators on $\CK $ with scalar product given by $\braket{a}{b}_2= \tr a^{\dagger}b$ for $a,b\in \Bt{\CK }$. If $\rho$ and $\sigma$ are density operators on $\CK $ with $\ker \sigma =\{0\}$ then the relative modular operator $\Delta_{\rho,\sigma}$ can be seen as the closure of densely defined  operator $X\sigma \mapsto \rho X$, $\forall X\in \B{\CK }$ and $\Delta_{\rho,\sigma}$ is a unbounded positive selfadjoint operator on the Hilbert space $\left(\Bt{\CK },  \braket{\cdot}{\cdot}_2\right)$ satisfying  \begin{align}\label{eq:relative-modular-action}
  \Delta_{\rho,\sigma} (X\sigma) = \rho X, \quad \forall X\in \B{\CK }.
\end{align}

Consider two density operators $\rho$ and $\sigma$ on $\CK $.  When considered as normal states $\varphi_{\rho}$ and $\varphi_{\sigma}$ on the von Neumann algebra $\B{\CK }$, let $E_{\rho,\sigma}$ denote the spectral measure associated with the relative modular operator $\Delta_{\rho,\sigma}$ of $\rho$ and $\sigma$. Araki defined the relative entropy \cite{araki1976relative,araki1977relative} of $\varphi_\rho$ and $\varphi_\sigma$  as 
\begin{align}
    \label{eq:araki-entropy}
S(\varphi_{\rho}||\varphi_\sigma):= \begin{cases}
 \int\limits_{0+}^{\infty} \log \lambda \mel{\sqrt{\rho}}{E_{\rho,\sigma}(\dd \lambda)}{\sqrt{\rho}}_2,& \text{ when } \supp \rho \subseteq \supp \sigma;\\
   \infty, &\text{ otherwise.} 
\end{cases}
\end{align}
 When $\CK $ is finite dimensional, it is well known \cite{araki1976relative,araki1977relative,petz1986quasientropy} that \begin{align}\label{eq:araki-equals-it=finite-dim}
     S(\varphi_{\rho}||\varphi_\sigma)= \begin{cases}
   \tr \rho^{1/2}(\log \rho) \rho^{1/2}-\tr\rho^{1/2}(\log \sigma)\rho^{1/2},&\text{ when } \supp\rho\subseteq\supp \sigma;\\
   \infty, &\text{ otherwise}.
 \end{cases}
 \end{align}
 Recently, there have been attempts \cite{luczak2019relative} to prove a similar formula as in (\ref{eq:araki-equals-it=finite-dim}) when $\dim \CK  =\infty$. In this section, we provide a counter example to show that (\ref{eq:araki-equals-it=finite-dim}) will not hold in general when $\CK $ is infinite dimensional. 


 
\begin{lem}\label{lem:relative-modular-spectral}
  Let $\rho$ and $\sigma$ be as in Example \ref{eg:counter-example}. Let $X_{1,j}=\ketbra{u_1}{v_j}\in \Bt{\CK }$, $j=1,2,\dots$. Then $\{X_{1,j}\}_{j=1}^{\infty}$ is an orthonormal set in $\Bt{\CK }$ and the spectral decomposition of $\Delta_{\rho,\sigma}$ is given by
  \begin{align}\label{eq:relative-modular-spectral}
   \Delta_{\rho,\sigma}= \sum_{j = 1}^{\infty} 2^j  \ketbra{X_{1,j}},
  \end{align}
  where $\ketbra{X_{1,j}}$ is  considered as a rank one projection acting on $\Bt{\CK }$. 
\end{lem}
 \begin{proof}
Extend $\{u_1\}$ to $\{u_1, u_2,\dots\}$ as an orthonormal basis of $\CK $. Define \begin{align}\label{eq:X-ell-j-defn}
    X_{\ell,j} = \ketbra{u_\ell}{v_j}, \quad \forall \ell, j = 1,2,\cdots.
\end{align}
Since $\{u_\ell\}$ and $\{v_j\}$ are orthonormal bases for $\CK $, it is easy to see that the double sequence \[\{X_{\ell,j}\}_{\ell, j=1}^{\infty} \text{ is an orthonormal basis of } \Bt{\CK }. \]
To complete the proof, we will show the following: \begin{enumerate}
    \item \label{item:counter-eg-1} $X_{\ell,j}\in D(\Delta_{\rho,\sigma})$ for all $\ell,j=1,2,\cdots$;
    \item \label{item:counter-eg-2} for $\ell\neq 1$, $X_{\ell,j}\in \ker \Delta_{\rho,\sigma}$;
    \item\label{item:counter-eg-3} $\Delta_{\rho,\sigma} (X_{1,j}) = 2^j X_{1,j}$ for all $j=1,2,\dots$.
\end{enumerate}


To prove \ref{item:counter-eg-1}, by (\ref{eq:relative-modular-action}) it is enough to prove that there exists $Y_{\ell,j}\in \B{\CK }$ such that $X_{\ell,j} = Y_{\ell,j}\sigma$. Note that \begin{align}\label{eq:X-ell-j}
    2^jX_{\ell,j}\sigma =2^j\ketbra{u_{\ell}}{v_j}\left(\sum_{k=1}^{\infty} 2^{-k}\ketbra{v_k}\right) =\ketbra{u_{\ell}}{v_j}= X_{\ell,j},\quad \forall \ell,j=1,2,\dots. 
\end{align}


Now to prove \ref{item:counter-eg-2} and \ref{item:counter-eg-3}, note that by  (\ref{eq:X-ell-j}) and (\ref{eq:relative-modular-action})
\begin{align*}
 \Delta_{\rho,\sigma} (X_{\ell,j})   &= \Delta_{\rho,\sigma} (2^jX_{\ell,j}\sigma)=2^j\rho X_{\ell,j}\\
 &=2^j\ketbra{u_1}\ketbra{u_{\ell}}{v_j}\\
 &=\begin{cases}
   2^j\ketbra{u_1}{v_j}, & \text{ if } \ell=1\\
   0, &\text{ if } \ell\neq 1.
 \end{cases}
\end{align*}
Thus $\Delta_{\rho,\sigma} (X_{1,j}) = 2^j X_{1,j}$ for all $j=1,2,\dots$ and $\Delta_{\rho,\sigma} (X_{\ell,j}) =0$ when $\ell\neq 1$.
 \end{proof}
 \begin{cor} Let $\rho$ and $\sigma$ be as in Example \ref{eg:counter-example} and $X_{\ell,j}$ be as in (\ref{eq:X-ell-j-defn}).
The spectral measure $E_{\rho,\sigma}$ of $\Delta_{\rho,\sigma}$ is supported on the set $\{0,2,2^2,\dots\}$ and \begin{align}
    \label{eq:spectral-measure-relative-modular}
    &E_{\rho,\sigma}\{0\} \phantom{.}=\sum_{\ell=2}^\infty\sum_{j=1}^{\infty}\ketbra{X_{\ell,j}},\\
     &E_{\rho,\sigma}\{2^j\}  =\ketbra{X_{1,j}} \quad j=1,2,\dots.
\end{align}
 \end{cor}
 \begin{proof}
 Follows easily from  Lemma \ref{lem:relative-modular-spectral} and the fact that $\{X_{\ell,j}\}_{\ell, j=1}^{\infty}$ forms an orthonormal basis of  $\Bt{\CK }$. 
 \end{proof}
 \begin{thm}\label{thm:counter-example}
Let $\rho$ and $\sigma$ be as in Example \ref{eg:counter-example}, then Araki's relative entropy of $\rho$ and $\sigma$ is  \[S(\varphi_{\rho}||\varphi_\sigma) = \infty,\] but
\[\tr \rho^{1/2}(\log \rho) \rho^{1/2}-\tr\rho^{1/2}(\log \sigma)\rho^{1/2}=0.\]
 \end{thm}
 \begin{proof} The second equality in the statement of the theorem is already shown in Example \ref{eg:counter-example}. To prove the first equality we will evaluate the integral in (\ref{eq:araki-entropy}). To this end, fix $j\in \BN$ and note that \begin{align*}
     \mel{\sqrt{\rho}}{E_{\rho,\sigma}\{2^j\}}{\sqrt{\rho}}_{2}&= \mel{\sqrt{\rho}}{\left(\ketbra{X_{1,j}}\right)}{\sqrt{\rho}}_{2}\\
     &=\braket{\sqrt{\rho}}{X_{1,j}}_2\braket{X_{1,j}}{\sqrt{\rho}}_2\\
     &=\abs{\braket{\sqrt{\rho}}{X_{1,j}}_2}^2\\
     &=\abs{\tr {\sqrt{\rho}}{X_{1,j}}}^2.
 \end{align*}
 Recall that $\rho=\ketbra{u_1}=\sqrt{\rho}$ and $X_{1,j}=\ketbra{u_1}{v_j}$. Therefore, \begin{align*}
     \mel{\sqrt{\rho}}{E_{\rho,\sigma}\{2^j\}}{\sqrt{\rho}}_{2} = \abs{\braket{u_1}{v_j}}^2.
 \end{align*}
Finally, since $u_1=\left(\sum_{k=1}^{\infty}\frac{1}{k^2}\right)^{-1/2}\sum_{k=1}^{\infty}\frac{1}{k}\ket{v_k}$, we have \begin{align}\label{eq:measure-araki-entropy}
   \mel{\sqrt{\rho}}{E_{\rho,\sigma}\{2^j\}}{\sqrt{\rho}}_{2} = \left(\sum_{k=1}^{\infty}\frac{1}{k^2}\right) \frac{1}{j^2}.
\end{align}
Now \begin{align*}
    S(\varphi_{\rho}||\varphi_\sigma)&=
 \int\limits_{0+}^{\infty} \log \lambda \mel{\sqrt{\rho}}{E_{\rho,\sigma}(\dd \lambda)}{\sqrt{\rho}}_2\\
 &=\sum_{j=1}^{\infty} \log 2^j\left(\sum_{k=1}^{\infty}\frac{1}{k^2}\right) \frac{1}{j^2}\\
 &=\left(\sum_{k=1}^{\infty}\frac{1}{k^2}\right)\sum_{j=1}^{\infty} \frac{j}{j^2}=\infty,
\end{align*}which completes the proof.
 \end{proof}
\subsubsection{Quantum Pinsker Type Inequality}
\begin{lem}\label{lem:total-variation}
 Let $\rho$ and $\sigma$ be as in (\ref{eq:spectral-rho-sigma})  and let $P$ and $Q$ be as in (\ref{eq:P-and-Q}), then the total variation distance $V(P,Q)$ between $P$ and $Q$ is given by  \begin{align}
     V(P,Q) = \sum_{i,j}\abs{r_i-s_j}\abs{\braket{u_i}{v_j}}^2
 \end{align}
\end{lem}
\begin{proof}
By the definition of the total variation distance as in  (\ref{eq:total-variation}), we have \begin{align*}
    V(P,Q) &= \sum_{x\in \mathcal{I}\times \mathcal{I}}\abs{\frac{P(x)}{\mu(x)}-\frac{Q(x)}{\mu(x)}}\mu(x) = \sum_{x\in \mathcal{I}\times \mathcal{I}}\abs{P(x)-Q(x)}\\
    &=\sum_{i,j}\abs{r_i-s_j}\abs{\braket{u_i}{v_j}}^2,
\end{align*}
which completes the proof.
\end{proof}
\begin{lem}\label{lem:tr-rho-sigma-square}
 Let $\rho$ and $\sigma$ be as in (\ref{eq:spectral-rho-sigma}), and let $P$ and $Q$ be as in (\ref{eq:P-and-Q}),  then \begin{align}
    \norm{\rho-\sigma}_2^2 = \sum_{i,j}(r_i-s_j)^2\abs{\braket{u_i}{v_j}}^2\leq V(P,Q),
 \end{align} 
  where $\norm{\cdot}_2$ denotes the Hilbert-Schmidt norm.
\end{lem}
\begin{proof}
Note that \begin{align}\label{eq:tr-rho-sigma-square}
    \norm{\rho-\sigma}_2^2 = \tr (\rho-\sigma)^2= \tr \rho^2+\tr\sigma^2 -2\tr \rho\sigma.
\end{align}
By (\ref{eq:spectral-rho-sigma}), \[\rho\sigma = \sum_{i,j}r_is_j\braket{u_i}{v_j}\ketbra{u_i}{v_j}\]
is a trace class operator. Hence \[\tr \rho\sigma = \sum_{i,j}r_is_j\abs{\braket{u_i}{v_j}}^2. \]
Along with the fact that $\tr\rho^2 = \sum_{i}r_i^2$ and $\tr \sigma^2 = \sum_j s_j^2$, we get from (\ref{eq:tr-rho-sigma-square}),\begin{align*}
    \norm{\rho-\sigma}_2^2 &=  \sum_{i}r_i^2+ \sum_j s_j^2-2\sum_{i,j}r_is_j\abs{\braket{u_i}{v_j}}^2\\
    &=\sum_{i,j}r_i^2\abs{\braket{u_i}{v_j}}^2+ \sum_{i,j} s_j^2\abs{\braket{u_i}{v_j}}^2-2\sum_{i,j}r_is_j\abs{\braket{u_i}{v_j}}^2\\
    &=\sum_{i,j} (r_i-s_j)^2\abs{\braket{u_i}{v_j}}^2.
\end{align*}
To prove the last inequality in (\ref{eq:tr-rho-sigma-square}) note that $r_i,s_j\in [0,1]$ for all $i,j$ and thus we have \[(r_i-s_j)^2\leq\abs{r_i-s_j}, \quad \forall i,j.\]  
Lemma \ref{lem:total-variation} completes the proof.
\end{proof}
\begin{thm}[Quantum Pinsker Type Inequality]\label{thm:quantum-pinsker}
 For $\alpha\in (0,1]$, the Petz-R\'enyi relative entropy satisfies the quantum Pinsker's inequality, \begin{align}
  \frac{\alpha}{2}\norm{\rho-\sigma}_2^4\leq   D_{\alpha}(\rho||\sigma),
 \end{align}
 where $\norm{\cdot}_2$ denotes the Hilbert-Schmidt norm.
\end{thm}
\begin{proof} Let $\rho$ and $\sigma$ be as in (\ref{eq:spectral-rho-sigma})  and let $P$ and $Q$ be as in (\ref{eq:P-and-Q}).
By Theorem \ref{thm:31-veh} and Theorem \ref{thm:quantum-reduces-to-classical} we have\[\frac{\alpha}{2}V^2(P,Q)\leq D_{\alpha}(P||Q) =D_{\alpha}(\rho||\sigma) ,\]
where $V(P,Q)$ denote the total variation distance between $P$ and $Q$.  Lemma \ref{lem:tr-rho-sigma-square} completes the proof.
\end{proof}
\begin{rmk}
The quantum Pinsker's inequality presented above  bounds the relative Petz-R\'{e}nyi entropy from
below by a function of the Hilbert-Schmidt norm of the difference of two quantum states. Usually, in the Pinsker's type inequalities the relative Petz-Renyi entropy between two states on a finite dimensional Hilbert space
is bounded from below by an increasing  function of their von Neumann relative entropy, which in turn, is bounded from below by a 
function of the trace class norm of their difference (for example one may refer to equation (43) in \cite{rastegin-2013}). Our inequality is weaker since we only use the Hilbert-Schmidt norm, but on the 
other hand, our inequality is obtained using only the classical Pinsker's inequality on the R\'{e}nyi relative entropy.
\end{rmk}

\subsubsection{Continuity, Positivity, Symmetry and Concavity} 
We begin this section with three examples which illustrate behaviour of $D_{\alpha}(\rho||\sigma)$ when $\alpha\geq 1.$ These examples will help us to understand continuity points of $D_{\alpha}(\rho||\alpha)$
\begin{eg}[$D_1(\rho||\sigma)<\infty$ but $D_{\alpha}(\rho||\sigma) = \infty, \forall \alpha >1$] Let $\{u_i\}_{i=1}^{\infty}$ be any orthonormal basis on $\CK $. Take \begin{align*}
    \rho &= \sum_{i=1}^{\infty}2^{-i}\ketbra{u_i}\\
    \sigma & =s^{-1} \sum_{j=1}^{\infty}2^{-j^2}\ketbra{u_j},
\end{align*}
 where $s=\left(\sum_j 2^{-j^2}\right)^{1/2}$. In this case, keeping the notations in (\ref{eq:spectral-rho-sigma}), and  (\ref{eq:P-and-Q}), $r_i=2^{-i}$, $s_j=s^{-1}2^{-j^2}$, $\braket{u_i}{v_j}= \delta_{i,j}$, we have
 \begin{align*}
     D_1(\rho||\sigma) = D(P||Q)&= \sum_{i}r_i\log \left( \frac{r_i}{s_i}\right) = \sum_{i}2^{-i}\log \left( \frac{s2^{-i}}{2^{-i^2}}\right)\\
     &= \sum_{i}2^{-i} \log \left(s2^{-i+i^2}\right)= \sum_{i}2^{-i} \log s+ \sum_{i}2^{-i} \log \left(2^{-i+i^2}\right) \\
     &= \sum_{i}2^{-i} \log s + \sum_{i}2^{-i}(-i+i^2) <\infty.
 \end{align*}
 On the other hand, for $\alpha>1$, \begin{align*}
     \sum_{i}r_i^{\alpha}s_i^{(1-\alpha)} &=s^{-1(1-\alpha)} \sum_{i}2^{-\alpha i}2^{-(1-\alpha)i^2} = s^{-1(1-\alpha)}\sum_{i}2^{(\alpha-1)i^2-\alpha i}=\infty
 \end{align*}
 because $(\alpha-1)>0$. Therefore,
 \begin{align*}
     D_{\alpha}(\rho||\sigma) = D_{\alpha}(P||Q)&= \frac{1}{\alpha-1} \log \sum_{i}r_i^{\alpha}s_i^{1-\alpha} = \infty. 
 \end{align*}
\end{eg}
\begin{eg}[$D_{\alpha}(\rho||\sigma)<\infty$ for $1<\alpha< 2$ but $D_{2}(\rho||\sigma) = \infty$] Let $\{u_i\}_{i=1}^{\infty}$ be any orthonormal basis on $\CK $. Take \begin{align*}
    \rho &= \sum_{i=1}^{\infty}2^{-i}\ketbra{u_i}\\
    \sigma & =s^{-1} \sum_{j=1}^{\infty}2^{-2i}\ketbra{u_j},
\end{align*}
 where $s=\left(\sum_j 2^{-2i}\right)^{1/2}$. In this case, keeping the notations in (\ref{eq:spectral-rho-sigma}), and  (\ref{eq:P-and-Q}), $r_i=2^{-i}$, $s_j=s^{-1}2^{-2i}$, $\braket{u_i}{v_j}= \delta_{i,j}$. We have
for $\alpha>1$, \begin{align*}
     \sum_{i}r_i^{\alpha}s_i^{(1-\alpha)} &=s^{-1(1-\alpha)} \sum_{i}2^{-\alpha i}2^{-(1-\alpha)2i} = s^{-1(1-\alpha)}\sum_{i}2^{(\alpha-1)2i-\alpha i}= s^{-1(1-\alpha)}\sum_{i}2^{(\alpha-2)i}.
 \end{align*}
 The above series converges for $1<\alpha<2$ and diverges for $\alpha=2$. Therefore,
 \begin{align*}
     D_{\alpha}(\rho||\sigma) = D_{\alpha}(P||Q)&= \frac{1}{\alpha-1} \log \sum_{i}r_i^{\alpha}s_i^{1-\alpha}
 \end{align*}
is finite for $1<\alpha<2$ and diverges for $\alpha=2$. 
\end{eg}
\begin{eg}[$D_{2}(\rho||\sigma)<\infty$  but $D_{\alpha}(\rho||\sigma) = \infty$ for $\alpha>2$] Let $\{u_i\}_{i=1}^{\infty}$ be any orthonormal basis on $\CK $. Take \begin{align*}
    \rho &= \sum_{i=1}^{\infty}2^{-i}\ketbra{u_i}\\
    \sigma & =s^{-1} \sum_{j=1}^{\infty}i^{2}2^{-2i}\ketbra{u_j},
\end{align*}
 where $s=\left(\sum_j i^{2}2^{-2i}\right)^{1/2}$. In this case, keeping the notations in (\ref{eq:spectral-rho-sigma}), and  (\ref{eq:P-and-Q}), $r_i=2^{-i}$, $s_j=s^{-1}i^{2}2^{-2i}$, $\braket{u_i}{v_j}= \delta_{i,j}$. We have
for $\alpha\geq 2$, \begin{align*}
     \sum_{i}r_i^{\alpha}s_i^{(1-\alpha)} &=s^{-1(1-\alpha)} \sum_{i}2^{-\alpha i}i^{2(1-\alpha)}2^{-(1-\alpha)2i}\\ 
     &= s^{-1(1-\alpha)}\sum_{i}i^{2(1-\alpha)}2^{(\alpha-1)2i-\alpha i}\\
     &= s^{-1(1-\alpha)}\sum_{i}i^{2(1-\alpha)}2^{(\alpha-2)i}.
 \end{align*}
 The above series converges for $\alpha=2$ and diverges for $\alpha>2$. Therefore,
 \begin{align*}
     D_{\alpha}(\rho||\sigma) = D_{\alpha}(P||Q)&= \frac{1}{\alpha-1} \log \sum_{i}r_i^{\alpha}s_i^{1-\alpha}
 \end{align*}
is finite for $\alpha=2$ and diverges for $\alpha>2$. 
\end{eg}
\begin{thm}
 \label{thm:7-veh-quantum}
%\label{thm:continuity-in-order}
The Petz-R\'enyi relative entropy $D_\alpha (\rho||\sigma)$ is continuous in $\alpha$ on $\mathcal{A} = \{\alpha\in [0,\infty]\,|\, 0\leq \alpha\leq 1\textnormal{ or } D_{\alpha}(\rho||\sigma)<\infty\}$.
\end{thm}
\begin{proof} By Theorem \ref{thm:quantum-reduces-to-classical} we know $D_\alpha (\rho||\sigma) = D_\alpha (P||Q)$ where $\rho$ and $\sigma$ are as in (\ref{eq:spectral-rho-sigma})  and $P$ and $Q$ are as in (\ref{eq:P-and-Q}). The result follows because the same result is true for the classical R\'enyi relative divergence (refer Theorem \ref{thm:5-veh}).
\end{proof}
\begin{thm}
\label{thm:8-veh-quantum} For any order $\alpha\in [0,\infty]$, 
\[D_{\alpha}(\rho||\sigma)\geq 0.\] For $\alpha>0$, $D_{\alpha}(\rho||\sigma) = 0$ if and only if $\rho=\sigma$. For $\alpha = 0$, $D_{\alpha}(\rho||\sigma) = 0$ if and only if $\supp \sigma \subseteq \supp \rho$.
\end{thm}
\begin{proof}
Follows easily from Theorem \ref{thm:8-veh} and Proposition \ref{prop:rho-equals-sigma-P-equals-Q} because of Theorem \ref{thm:quantum-reduces-to-classical}.
\end{proof}
\begin{prop}
\label{prop:8-veh-quantum}
For any $0<\alpha<1$, the Petz-R\'enyi relative entropy shows the following skew-symmetry property \[D_\alpha(\rho||\sigma) = \frac{\alpha}{1-\alpha}D_{1-\alpha}(\sigma||\rho).\]
\end{prop}
\begin{proof}
Follows from Proposition \ref{prop:8-veh}.
\end{proof}
Note that in particular, Petz-R\'enyi relative entropy is symmetric for $\alpha=1/2$, and that skew-symmetry does not hold for $\alpha = 0$ and $\alpha = 1$. 
\begin{thm}
\label{thm:16-veh-quantum}
For any $0<\alpha\leq \beta <1$,\[\frac{\alpha}{\beta}\frac{1-\beta}{1-\alpha}D_{\beta}(\rho||\sigma)\leq D_\alpha (\rho||\sigma)\leq D_{\beta}(\rho||\sigma).\]
\end{thm}
\begin{proof}
Follows from Theorems \ref{thm:16-veh} and \ref{thm:quantum-reduces-to-classical}.
\end{proof}
\begin{rmk}In the light of the Theorem~\ref{thm:16-veh-quantum} we can discuss about a topology on the set of states arising from $D_{\alpha}$.
For a fixed $\alpha \in (0,1)$, one can define $\alpha$-left open ball with center $\rho$ and radius $r>0$ to be the set
$\{ \sigma\,|\, D_\alpha(\rho||\sigma) <r\}$, and subsequently define $\alpha$-left open sets to be the union of $\alpha$-left open balls.
Notice that Theorem~\ref{thm:16-veh-quantum} yields that for $\alpha, \beta \in (0,1)$, the $\alpha$-left topology is equivalent to the 
$\beta$-left topology. Similarly, one can define $\alpha$-right open balls and $\alpha$-right topologies by reversing the order 
of $\rho$ and $\sigma$ in the definition of $\alpha$-left topology. Proposition~\ref{prop:8-veh-quantum}, combined with the fact that the 
$\alpha$-left topologies are all equivalent for $0<\alpha<1$, gives that the $\alpha$-left topologies are equivalent with the 
$\beta$-right topologies for all $\alpha, \beta \in (0,1)$.
\end{rmk}
\begin{thm}
\label{thm:23-veh-quantum}The following conditions are equivalent:
\begin{enumerate}
    \item \label{item:vEH-1-quantum}$\supp \sigma \subseteq \supp \rho,$
    \item \label{item:vEH-2-quantum} $\tr \Pi_\rho\sigma = 1$, where $\Pi_\rho$ is the orthogonal projection onto $\supp \rho$.
    \item\label{item:vEH-3-quantum} $D_0(\rho||\sigma) = 0$,
    \item\label{item:vEH-4-quantum} $\lim_{\alpha \downarrow 0} D_{\alpha}(\rho||\sigma) = 0.$
\end{enumerate}
\end{thm}
\begin{proof}
\ref{item:vEH-1-quantum} $\Leftrightarrow$ \ref{item:vEH-3-quantum} follows from Theorem \ref{thm:8-veh-quantum}. 

\ref{item:vEH-3-quantum} $\Leftrightarrow$ \ref{item:vEH-2-quantum} follows from Theorem \ref{thm:4-veh-quantum}. 


\ref{item:vEH-3-quantum} $\Leftrightarrow$ \ref{item:vEH-4-quantum} follows from Theorem \ref{thm:7-veh-quantum}.
\end{proof}
\begin{thm}
\label{thm:24-veh-quantum}The following conditions are equivalent:
\begin{enumerate}
    \item \label{item:vEH-5-quantum}$\supp \rho\perp \supp \sigma,$
    \item \label{item:vEH-6-quantum} $\tr \Pi_\rho\sigma = 0$, where $\Pi_\rho$ is the orthogonal projection onto $\supp \rho$,
    \item\label{item:vEH-7-quantum} $D_{\alpha}(\rho||\sigma) = \infty$ for some $\alpha\in [0,1)$,
    \item\label{item:vEH-8-quantum} $D_{\alpha}(\rho||\sigma) = \infty$ for all $\alpha\in [0,\infty].$
\end{enumerate}
\end{thm}
\begin{proof}
\ref{item:vEH-5-quantum}  $\Leftrightarrow$ $\supp \sigma \subseteq \ran (I-\Pi_{\rho})$ $\Leftrightarrow$ $\tr (I-\Pi_{\rho})\sigma = \tr \sigma =1$ $\Leftrightarrow$ the statement  \ref{item:vEH-6-quantum}.

By Theorem \ref{thm:4-veh-quantum}, the statement \ref{item:vEH-6-quantum} $\Leftrightarrow$ $D_0(\rho||\sigma) = \infty$  $\Leftrightarrow$ statement \ref{item:vEH-8-quantum} by Theorem \ref{thm:3-veh-quantum}.

Finally, \ref{item:vEH-7-quantum} $\Leftrightarrow$ \ref{item:vEH-8-quantum} because of Theorem \ref{thm:16-veh-quantum}


\end{proof}
\begin{prop}\label{cor:2-veh-quantum}
The function $[0, \infty] \ni \alpha \mapsto (1-\alpha)D_{\alpha}(\rho||\sigma)$ is concave, with the conventions that it is $0$ at $\alpha = 1$ even if $D(\rho||\sigma) = \infty$ and that it is $0$ at $\alpha = \infty$ if $\rho=\sigma$.

\end{prop}
\begin{proof}
Follows from Corollary \ref{cor:2-veh} and Theorem \ref{thm:quantum-reduces-to-classical}.
\end{proof}
\section{Conclusion and Discussion}
The major contributions of this article are: (1) we provide a procedure for obtaining results related to Petz-R\'enyi relative entropy using the classical results available for R\'enyi divergence of probability measures; and (2) we give an example to show  that Araki's definition of von Neumann relative entropy in infinite dimensions does not agree with its usual information theoretic definition.
We illustrated (1)  by proving several results about Petz-R\'enyi relative entropy from the classical results available in the survey article \cite{Erven-Harremos-2014}. All our results work both in finite and infinite dimensional setting which is a strength of this method. Hence these results are particularly useful in continuous variable quantum information theory as well. For instance, we use the results of the current work in order to study the Petz-R\'enyi relative entropy of gaussian states in a follow-up article \cite{androulakis-john-2022}.  There have been several recent new results in the classical theory of R\'enyi divergences after the publication of \cite{Erven-Harremos-2014}. This shows that more results about Petz-R\'enyi relative entropy can be proved using the Theorem \ref{thm:quantum-reduces-to-classical} of this article. This method directly works for proving a quantum counterpart of any result about the classical R\'enyi divergence involving only two probability distributions. An interesting open problem is to develop methods in order to overcome the technical difficulties one encounters while trying to obtain a quantum version of a result  starting from a classical result in which more than two probability distributions are involved. One example of such  result, that may be explored is the data processing inequality which is satisfied by many other entropic quantities.
\section*{Acknowledgement}
The second author thanks The Fulbright Scholar Program and United States-India Educational Foundation for providing funding and other support to conduct this research through a Fulbright-Nehru Postdoctoral Fellowship.







\appendix
\section{Appendix}
\subsection{R\'enyi divergence  in the Classical setting}\label{sec:classical-renyi-devergence}
In this section, we recall a few facts about the R\'enyi divergence in the setting of classical probability. We refer to the survey article \cite{Erven-Harremos-2014} for the following definitions and results which we state in this subsection. The results we need from \cite{Erven-Harremos-2014}  are repeated here for the ease of the reader. If $\mu$ and $\nu$ are two positive measures on a measure space $(X,\mathcal{F})$, then $\nu$ is said to be absolutely continuous with respect to $
\mu$ and we write $\nu\ll \mu$, if for every $E\in \mathcal{F}$ such that $\mu (E) = 0$, then $\nu(E) = 0$.
\begin{defn}\label{defn:renyi-divergence-Kullback}
 Let $P, Q$ be probability distributions on a measure space $(X, \mathcal{F})$. Let $\mu$ be any $\sigma$-finite measure such that $P\ll\mu$ and $Q\ll\mu$. Let $p$ and $q$  denote the Radon-Nikodym derivatives with respect to $\mu$, of $P$ and $Q$, respectively. \begin{enumerate}
     \item The R\'enyi divergence of order $\alpha\in (0,1)\cup (1,\infty)$ is defined as \begin{equation}
     \label{eq:classical-renyi-alpha}
     D_{\alpha}(P||Q) = \frac{1}{\alpha-1} \log \int_X p^\alpha q^{1-\alpha} \dd \mu,
 \end{equation}
 where for $\alpha>1$, we adopt the conventions $0^{1-\alpha} = \infty$ and $0\cdot \infty = 0$.
 \item\label{defn:kullback-leibler} The Kullback-Leibler divergence of $P$ from $Q$ is defined as 
 \begin{equation}
  D(P||Q) = \int p\log \frac{p}{q} \dd \mu,
 \end{equation}
 with the conventions that $0\log (0/q) = 0$, for $q\geq0$ and $p\log(p/0)=\infty$ if $p>0$. Consequently, \begin{align}\label{eq:kullback-leibler-equals-infty}
     D(P||Q) = \infty \textnormal{ if } P\nll Q.
 \end{align}
 \end{enumerate}    
\end{defn}
\begin{rmk}
The R\'enyi divergence defined by (\ref{eq:classical-renyi-alpha}) is independent of the choice of $\mu$. To see this, fix the measure $\mu_0=\frac{P+Q}{2}$, then $\mu_0$ is a probability measure on $(X, \mathcal{F})$ such that $P\ll\mu_0$ and  $Q\ll \mu_0.$ 
We will show that $\int \left(\dv{P}{\mu}\right)^{\alpha}\left(\dv{Q}{\mu}\right)^{1-\alpha}\dd \mu$ $=$ $\int \left(\dv{P}{\mu_0}\right)^{\alpha}\left(\dv{Q}{\mu_0}\right)^{1-\alpha}\dd \mu_0$ for any measure $\mu$ such that $P\ll\mu$ and $Q\ll \mu$. Since $P\ll\mu$ and $Q\ll \mu$, we have $P+Q\ll \mu$. Hence $\mu_0\ll\mu$ and we have 
\begin{align*}
\int \left(\dv{P}{\mu_0}\right)^{\alpha}\left(\dv{Q}{\mu_0}\right)^{1-\alpha}\dd \mu_0 & =   \int \left(\dv{P}{\mu_0}\right)^{\alpha}\left(\dv{Q}{\mu_0}\right)^{1-\alpha}\left(\dv{\mu_0}{\mu}\right)\dd \mu\\
&=\int \left(\dv{P}{\mu_0}\dv{\mu_0}{\mu}\right)^{\alpha}\left(\dv{Q}{\mu_0}\dv{\mu_0}{\mu}\right)^{1-\alpha}\dd \mu\\
&=\int \left(\dv{P}{\mu}\right)^{\alpha}\left(\dv{Q}{\mu}\right)^{1-\alpha}\dd \mu.
\end{align*}
\end{rmk}
\begin{defn}\label{defn:extended-orders}
 The R\'enyi divergences of orders $0, 1$ and $\infty$ are defined as \begin{align*}
     D_0(P||Q) &= \lim\limits_{\alpha \downarrow 0} D_{\alpha}(P||Q)\\
     D_1(P||Q) &= \lim\limits_{\alpha \uparrow 1} D_{\alpha}(P||Q)\\
     D_{\infty}(P||Q) &= \lim\limits_{\alpha \uparrow \infty} D_{\alpha}(P||Q)
 \end{align*}
\end{defn}
The limits in Definition \ref{defn:extended-orders} always exist because R\'enyi divergence is nondecreasing in order. 
\begin{thm}[Theorem 3 in \cite{Erven-Harremos-2014}] \label{thm:3-veh}
For $\alpha\in [0, \infty]$ the R\'enyi divergence $D_{\alpha}(P||Q)$ is nondecreasing in $\alpha$.
\end{thm}
\begin{thm}
[Theorem 4 in \cite{Erven-Harremos-2014}] \label{thm:4-veh}
\begin{equation}\label{eq:limit-at-0}D_0 (P||Q) = -\log (Q ({\{p>0\}})).
\end{equation}
\end{thm}   
\begin{thm}[Theorem 5 in \cite{Erven-Harremos-2014}]\label{thm:5-veh}
The Kullback-Leibler divergence is the limit of the R\'enyi divergence, i.e.,
\begin{equation}\label{eq:kullback-equals-left-limit-1}
    D(P||Q)= D_1(P||Q).
\end{equation}
Moreover, if $D(P||Q) = \infty$ or there exists $\beta>1$ such that $D_\beta (P||Q)<\infty$, then also \begin{equation}\label{eq:kullback-equals-right-limit-1}
    \lim\limits_{\alpha \downarrow 1} D_{\alpha}(P||Q) = D(P||Q).
\end{equation}
\end{thm}
\begin{rmk}
It is possible that $D_{\alpha}(P||Q) = \infty$ for all $\alpha>1$, but $D(P||Q)< \infty$, and hence (\ref{eq:kullback-equals-right-limit-1}) does not hold \cite{Erven-Harremos-2014}.
\end{rmk}
For any random variable $Y$, the \emph{essential supremum of} $Y$ with respect to $P$ is $\underset{P}{\operatorname{ess\ sup}}  Y = \sup\{c|P(Y>c)>0\}$.
\begin{thm}
[Theorem 6 in \cite{Erven-Harremos-2014}]\label{thm:6-veh}
\begin{equation}
 \label{eq:limit-infty}
 D_{\infty}(P||Q) = \log \underset{A\in \mathcal{F}}{\sup}\frac{P(A)}{Q(A)} = \log \left(\underset{P}{\operatorname{ess\ sup}} \frac{p}{q}\right),
\end{equation}
with the convention that $0/0 = 0$ and $x/0 = \infty$ if $x>0$.
\end{thm}
If the sample space $X$ is countable, then with the notational conventions of  Theorem~\ref{thm:6-veh} the 
$P$-essential supremum of $\frac{p}{q}$ reduces to the ordinary supremum of $\frac{p}{q}$, which in turn is equal to the 
supremum of $\frac{P}{Q}$, and we have 
\begin{equation}\label{eq:limit-infty-countable-sample-space}
    D_{\infty}(P||Q) = \log   \sup_x \frac{P(x)}{Q(x)},
\end{equation}
with the convention that $0/0 = 0$ and $x/0 = \infty$ if $x>0$.
\begin{thm}
[Theorem 7 in \cite{Erven-Harremos-2014}]\label{thm:7-veh}
The R\'enyi divergence $D_\alpha (P||Q)$ is continuous in $\alpha$ on $\mathcal{A} = \{\alpha\in [0,\infty]\,|\, 0\leq \alpha\leq 1\textnormal{ or } D_{\alpha}(P||Q)<\infty\}$.
\end{thm}
\begin{thm}
[Theorem 8 in \cite{Erven-Harremos-2014}]\label{thm:8-veh} For any order $\alpha\in [0,\infty]$ 
\[D_{\alpha}(P||Q)\geq 0.\] For $\alpha>0$, $D_{\alpha}(P||Q) = 0$ if and only if $P=Q$. For $\alpha = 0$, $D_{\alpha}(P||Q) = 0$ if and only if $Q\ll P$.
\end{thm}
\begin{prop}
[Proposition 2 in \cite{Erven-Harremos-2014}]\label{prop:8-veh}
For any $0<\alpha<1$, the R\'enyi divergence shows the following skew-symmetry property \[D_\alpha(P||Q) = \frac{\alpha}{1-\alpha}D_{1-\alpha}(Q||P).\]
\end{prop}
Note that in particular, R\'enyi divergence is symmetric for $\alpha=1/2$, and that skew-symmetry does not hold for $\alpha = 0$ and $\alpha = 1$. 


\begin{thm}
[Theorem 16 in \cite{Erven-Harremos-2014}]\label{thm:16-veh}
For any $0<\alpha\leq \beta <1$,\[\frac{\alpha}{\beta}\frac{1-\beta}{1-\alpha}D_{\beta}(P||Q)\leq D_\alpha (P||Q)\leq D_{\beta}(P||Q).\]
\end{thm}

\begin{thm}
[Theorem 23 in \cite{Erven-Harremos-2014}]\label{thm:23-veh}The following conditions are equivalent:
\begin{enumerate}
    \item \label{item:vEH-1}$Q\ll P,$
    \item \label{item:vEH-2} $Q(\{p>0\}) = 1$,
    \item\label{item:vEH-3} $D_0(P||Q) = 0$,
    \item\label{item:vEH-4} $\lim_{\alpha \downarrow 0} D_{\alpha}(P||Q) = 0.$
\end{enumerate}
\end{thm}
\begin{thm}
[Theorem 24 in \cite{Erven-Harremos-2014}]\label{thm:24-veh}The following conditions are equivalent:
\begin{enumerate}
    \item \label{item:vEH-5}$P\perp Q,$
    \item \label{item:vEH-6} $Q(\{p>0\}) = 0$,
    \item\label{item:vEH-7} $D_{\alpha}(P||Q) = \infty$ for some $\alpha\in [0,1)$,
    \item\label{item:vEH-8} $D_{\alpha}(P||Q) = \infty$ for all $\alpha\in [0,\infty].$
\end{enumerate}
\end{thm}
\begin{cor}[Corollary 2 in \cite{Erven-Harremos-2014}]\label{cor:2-veh}
The function $[0, \infty] \ni \alpha \mapsto (1-\alpha)D_{\alpha}(P||Q)$ is concave, with the conventions that it is $0$ at $\alpha = 1$ even if $D(P||Q) = \infty$ and that it is $0$ at $\alpha = \infty$ if $P=Q$.
\end{cor}
Recall the \emph{total variation distance} between $P$ and $Q$, defined as \begin{equation}\label{eq:total-variation}
    V(P,Q) = \int \abs{p-q} \dd \mu = 2\underset{A\in \mathcal{F}}{\sup} \abs{P(A)-Q(A)}.
\end{equation}
With this definition of total variation distance we now state the theorem on \emph{Pinsker's Inequality}.
\begin{thm}
[Theorem 31 in \cite{Erven-Harremos-2014}]\label{thm:31-veh}
Let $V(P,Q)$ be the total variation distance as in (\ref{eq:total-variation}). Then for any $\alpha \in (0,1],$
\begin{equation}
    \frac{\alpha}{2}V^2(P,Q) \leq D_{\alpha}(P||Q).
\end{equation}
\end{thm}
\subsection{Relative Modular Operator on \texorpdfstring{\ensuremath{\B{\CK }}}{}}\label{appendix:relative-modular}Let $\rho$ and $\sigma$ be two density operators  on a Hilbert space $\CK $. Assume further that $\ker \sigma =\{0\}$.
In this section, we provide a quick exposition to the definition of the relative modular $\Delta_{\rho,\sigma}$. The references for this exposition are \cite{araki1977relative, luczak2019relative}. Let $(\Bt{\CK }, \braket{\cdot}{\cdot}_2)$ denote the Hilbert space of Hilbert-Schmidt operators on $\CK $. Then  define the closable antilinear operator $S$ on the dense domain \[D(S)= \{X\sqrt{\sigma}|X\in \B{\CK }\}\subseteq \Bt{\CK }\]
by \begin{align}\label{eq:antilinear-S}
    S(X\sqrt{\sigma}) =X^{\dagger}\sqrt{\rho}.
\end{align}
For an antilinear operator $S$, the adjoint $S^{\dagger}$ is defined as the operator satisfying \[\braket{a}{Sb} = \overline{\braket{S^{\dagger}a}{b}}\]
on the appropriate domains. Now for the the antilinear operator $S$ defined by (\ref{eq:antilinear-S}) it may be noticed that \[\braket{\sqrt{\sigma}Y}{S(X\sqrt{\sigma})}_2 = \tr (Y^{\dagger}\sqrt{\sigma}X^{\dagger}\sqrt{\rho}) = \overline{\braket{\sqrt{\rho}Y^\dagger}{X\sqrt{\sigma}}_2}, \quad \forall X,Y\in \B{\CK }.\]
Therefore, $S^{\dagger}$ is also densely defined and satisfies 
\begin{align}\label{eq:antilinear-S-dagger}
    S^{\dagger}\left(\sqrt{\sigma}Y\right) =\sqrt{\rho}Y^{\dagger}.
\end{align}
The relative modular operator $\Delta_{\rho,\sigma}$ is defined as \begin{align}
    \label{eq:modular-operator}
    \Delta_{\rho,\sigma} =S^{\dagger}\overline{S}
\end{align}
where $\overline{S}$ denotes the closure of $S$. Furthermore, it may also be noted that 
\begin{align*}\label{eq:modular-operator-action}
    \Delta_{\rho,\sigma} \left(X\sigma\right) &= S^{\dagger}\overline{S}\left(X\sqrt{\sigma}\sqrt{\sigma}\right)\\
    &= S^{\dagger}\left(\sqrt{\sigma}X^{\dagger}\sqrt{\rho}\right)\\
    &= \sqrt{\rho}\left(X^{\dagger}\sqrt{\rho}\right)^\dagger\\
    &=\rho X, \quad \forall X\in \B{\CK }\numberthis.
\end{align*}
\bibliographystyle{IEEEtran}
\bibliography{bibliography}
\end{document}

%%% Local Variables:
%%% mode: latex
%%% TeX-master: t
%%% End:

% Generated by IEEEtran.bst, version: 1.14 (2015/08/26)
\begin{thebibliography}{10}
\providecommand{\url}[1]{#1}
\csname url@samestyle\endcsname
\providecommand{\newblock}{\relax}
\providecommand{\bibinfo}[2]{#2}
\providecommand{\BIBentrySTDinterwordspacing}{\spaceskip=0pt\relax}
\providecommand{\BIBentryALTinterwordstretchfactor}{4}
\providecommand{\BIBentryALTinterwordspacing}{\spaceskip=\fontdimen2\font plus
\BIBentryALTinterwordstretchfactor\fontdimen3\font minus
  \fontdimen4\font\relax}
\providecommand{\BIBforeignlanguage}[2]{{%
\expandafter\ifx\csname l@#1\endcsname\relax
\typeout{** WARNING: IEEEtran.bst: No hyphenation pattern has been}%
\typeout{** loaded for the language `#1'. Using the pattern for}%
\typeout{** the default language instead.}%
\else
\language=\csname l@#1\endcsname
\fi
#2}}
\providecommand{\BIBdecl}{\relax}
\BIBdecl

\bibitem{kullback-leibler-1951}
\BIBentryALTinterwordspacing
S.~Kullback and R.~A. Leibler, ``On information and sufficiency,'' \emph{Ann.
  Math. Statistics}, vol.~22, pp. 79--86, 1951. [Online]. Available:
  \url{https://doi.org/10.1214/aoms/1177729694}
\BIBentrySTDinterwordspacing

\bibitem{renyi1961}
A.~R\'{e}nyi, ``On measures of entropy and information,'' in \emph{Proc. 4th
  {B}erkeley {S}ympos. {M}ath. {S}tatist. and {P}rob., {V}ol. {I}}.\hskip 1em
  plus 0.5em minus 0.4em\relax Univ. California Press, Berkeley, Calif., 1961,
  pp. 547--561.

\bibitem{umegaki-1962}
\BIBentryALTinterwordspacing
H.~Umegaki, ``Conditional expectation in an operator algebra. {IV}. {E}ntropy
  and information,'' \emph{Kodai Math. Sem. Rep.}, vol.~14, pp. 59--85, 1962.
  [Online]. Available: \url{http://projecteuclid.org/euclid.kmj/1138844604}
\BIBentrySTDinterwordspacing

\bibitem{araki1976relative}
H.~Araki, ``Relative entropy of states of von neumann algebras,''
  \emph{Publications of the Research Institute for Mathematical Sciences},
  vol.~11, no.~3, pp. 809--833, 1976.

\bibitem{araki1977relative}
------, ``Relative entropy for states of von neumann algebras ii,''
  \emph{Publications of the Research Institute for Mathematical Sciences},
  vol.~13, no.~1, pp. 173--192, 1977.

\bibitem{petz1986quasientropy}
\BIBentryALTinterwordspacing
D.~Petz, ``Quasi-entropies for finite quantum systems,'' \emph{Reports on
  Mathematical Physics}, vol.~23, no.~1, pp. 57--65, 1986. [Online]. Available:
  \url{https://www.sciencedirect.com/science/article/pii/0034487786900674}
\BIBentrySTDinterwordspacing

\bibitem{petz-ohya-1993}
\BIBentryALTinterwordspacing
M.~Ohya and D.~Petz, \emph{Quantum entropy and its use}, ser. Texts and
  Monographs in Physics.\hskip 1em plus 0.5em minus 0.4em\relax
  Springer-Verlag, Berlin, 1993. [Online]. Available:
  \url{https://doi.org/10.1007/978-3-642-57997-4}
\BIBentrySTDinterwordspacing

\bibitem{Datta2009}
\BIBentryALTinterwordspacing
N.~Datta, ``Min- and max-relative entropies and a new entanglement monotone,''
  \emph{IEEE Trans. Inf. Theor.}, vol.~55, no.~6, p. 2816–2826, jun 2009.
  [Online]. Available: \url{https://doi.org/10.1109/TIT.2009.2018325}
\BIBentrySTDinterwordspacing

\bibitem{Ses-Lam-Wil-2018}
\BIBentryALTinterwordspacing
K.~P. Seshadreesan, L.~Lami, and M.~M. Wilde, ``Rényi relative entropies of
  quantum gaussian states,'' \emph{Journal of Mathematical Physics}, vol.~59,
  no.~7, p. 072204, 2018. [Online]. Available:
  \url{https://doi.org/10.1063/1.5007167}
\BIBentrySTDinterwordspacing

\bibitem{Par2021a}
K.~R. {Parthasarathy}, ``{A pedagogical note on the computation of relative
  entropy of two $n$-mode gaussian states},'' \emph{arXiv e-prints}, p.
  arXiv:2102.06708, Feb. 2021.

\bibitem{muller-tomamichel-etal-2013}
\BIBentryALTinterwordspacing
M.~M\"{u}ller-Lennert, F.~Dupuis, O.~Szehr, S.~Fehr, and M.~Tomamichel, ``On
  quantum {R}\'{e}nyi entropies: a new generalization and some properties,''
  \emph{J. Math. Phys.}, vol.~54, no.~12, pp. 122\,203, 20, 2013. [Online].
  Available: \url{https://doi.org/10.1063/1.4838856}
\BIBentrySTDinterwordspacing

\bibitem{frank-leib-2013}
\BIBentryALTinterwordspacing
R.~L. Frank and E.~H. Lieb, ``Monotonicity of a relative rényi entropy,''
  \emph{Journal of Mathematical Physics}, vol.~54, no.~12, p. 122201, 2013.
  [Online]. Available: \url{https://doi.org/10.1063/1.4838835}
\BIBentrySTDinterwordspacing

\bibitem{wilde-winter-yang-2014}
\BIBentryALTinterwordspacing
M.~M. Wilde, A.~Winter, and D.~Yang, ``Strong converse for the classical
  capacity of entanglement-breaking and {H}adamard channels via a sandwiched
  {R}\'{e}nyi relative entropy,'' \emph{Comm. Math. Phys.}, vol. 331, no.~2,
  pp. 593--622, 2014. [Online]. Available:
  \url{https://doi.org/10.1007/s00220-014-2122-x}
\BIBentrySTDinterwordspacing

\bibitem{Datta_Leditzky_2014}
\BIBentryALTinterwordspacing
N.~Datta and F.~Leditzky, ``A limit of the quantum r{\'{e}}nyi divergence,''
  \emph{Journal of Physics A: Mathematical and Theoretical}, vol.~47, no.~4, p.
  045304, jan 2014. [Online]. Available:
  \url{https://doi.org/10.1088/1751-8113/47/4/045304}
\BIBentrySTDinterwordspacing

\bibitem{Par2021b}
K.~R. {Parthasarathy}, ``{Computation of sandwiched relative alpha-entropy of
  two n-mode gaussian states},'' \emph{arXiv e-prints}, p. arXiv:2109.04062,
  Sep. 2021.

\bibitem{petz-hiai-1991}
\BIBentryALTinterwordspacing
F.~Hiai and D.~Petz, ``The proper formula for relative entropy and its
  asymptotics in quantum probability,'' \emph{Comm. Math. Phys.}, vol. 143,
  no.~1, pp. 99--114, 1991. [Online]. Available:
  \url{http://projecteuclid.org/euclid.cmp/1104248844}
\BIBentrySTDinterwordspacing

\bibitem{luczak2019relative}
A.~{\L}uczak, H.~Pods{\k{e}}dkowska, and R.~Wieczorek, ``Relative modular
  operator in semifinite von neumann algebras and its use,'' \emph{arXiv
  preprint arXiv:1912.09633}, 2019.

\bibitem{vN2018}
J.~von Neumann, \emph{Mathematical Foundations of Quantum Mechanics}, new
  edition~ed.\hskip 1em plus 0.5em minus 0.4em\relax Princeton University
  Press, 2018.

\bibitem{pedersen2012analysis}
\BIBentryALTinterwordspacing
G.~Pedersen, \emph{Analysis Now}, ser. Graduate Texts in Mathematics.\hskip 1em
  plus 0.5em minus 0.4em\relax Springer New York, 2012. [Online]. Available:
  \url{https://books.google.com/books?id=zNbiBwAAQBAJ}
\BIBentrySTDinterwordspacing

\bibitem{ReS75ii}
\BIBentryALTinterwordspacing
M.~Reed and B.~Simon, \emph{{I}{I}: Fourier Analysis, Self-Adjointness}, ser.
  Methods of Modern Mathematical Physics.\hskip 1em plus 0.5em minus
  0.4em\relax Elsevier Science, 1975. [Online]. Available:
  \url{https://books.google.com/books?id=Kz7s7bgVe8gC}
\BIBentrySTDinterwordspacing

\bibitem{Par12}
K.~R. Parthasarathy, \emph{An introduction to quantum stochastic calculus},
  ser. Modern Birkh\"auser Classics.\hskip 1em plus 0.5em minus 0.4em\relax
  Birkh\"auser/Springer Basel AG, Basel, 1992, [2012 reprint of the 1992
  original] [MR1164866].

\bibitem{rastegin-2013}
\BIBentryALTinterwordspacing
A.~E. Rastegin, ``Bounds of the {P}insker and {F}annes types on the {T}sallis
  relative entropy,'' \emph{Math. Phys. Anal. Geom.}, vol.~16, no.~3, pp.
  213--228, 2013. [Online]. Available:
  \url{https://doi.org/10.1007/s11040-013-9128-z}
\BIBentrySTDinterwordspacing

\bibitem{Erven-Harremos-2014}
T.~van Erven and P.~Harremos, ``Rényi divergence and kullback-leibler
  divergence,'' \emph{IEEE Transactions on Information Theory}, vol.~60, no.~7,
  pp. 3797--3820, 2014.

\bibitem{androulakis-john-2022}
G.~Androulakis and T.~C. John, ``Petz-r\'enyi relative entropy of quantum
  gaussian states,'' \emph{In Preparation}, 2022.

\end{thebibliography}
